\documentclass[12pt,a4paper,twoside]{article}

% --- Encoding & Sprache ---
\usepackage[utf8]{inputenc}
\usepackage[T1]{fontenc}
\usepackage[ngerman]{babel}
\usepackage{csquotes}

% --- Schrift & Layout ---
\usepackage{lmodern}
\usepackage{microtype}
\usepackage[margin=2.5cm]{geometry}
\usepackage{setspace}

% --- Mathematik ---
\usepackage{amsmath, amssymb, amsthm}

% --- Grafiken & Tabellen ---
\usepackage{graphicx}
\usepackage{booktabs}
\usepackage{longtable}
\usepackage{tabularx}
\usepackage{multirow}
\usepackage{array}
\usepackage{microtype}
\usepackage{xcolor}
\usepackage{colortbl}
\definecolor{autosarblue}{RGB}{0,83,159}
\definecolor{pythonblue}{RGB}{53,114,165}
\definecolor{lightgray}{RGB}{240,240,240}

% --- Code-Listing ---
\usepackage{listings}
\lstset{
  basicstyle=\ttfamily\footnotesize,
  breaklines=true,
  frame=single,
  numbers=left,
  numberstyle=\tiny\color{gray},
  keywordstyle=\color{autosarblue}\bfseries,
  commentstyle=\color{gray}\itshape,
  stringstyle=\color{pythonblue},
  showstringspaces=false,
  tabsize=4,
  captionpos=b,
  backgroundcolor=\color{lightgray}
}
\lstdefinelanguage{MicroPython}{
  language=Python,
  morekeywords={micropython, const, native, viper, machine, Pin, UART,
                SPI, I2C, Timer, RTC, WDT, ADC, DAC, PWM},
  sensitive=true,
}

% --- Hyperlinks & PDF ---
\usepackage[
  pdftitle={Python und MicroPython als Standardsprache in AUTOSAR},
  pdfauthor={Wissenschaftliche Arbeit},
  pdfsubject={AUTOSAR, Python, MicroPython, Embedded Systems, Standardisierung},
  pdfkeywords={AUTOSAR, Python, MicroPython, Embedded, Automotive, Safety, Standardisierung},
  colorlinks=true,
  linkcolor=autosarblue,
  citecolor=autosarblue,
  urlcolor=pythonblue
]{hyperref}

% --- Literaturverzeichnis ---
\usepackage{cite}

% --- DOI-Befehl ---
\newcommand{\doi}[1]{\url{https://doi.org/#1}}

% --- Theorem-Umgebungen ---
\newtheorem{definition}{Definition}[section]
\newtheorem{theorem}{Satz}[section]
\newtheorem{proposition}{Proposition}[section]
\newtheorem{remark}{Bemerkung}[section]

% --- Abkürzungsverzeichnis ---
\usepackage[acronym,toc]{glossaries}
\makenoidxglossaries
\newacronym{autosar}{AUTOSAR}{AUTomotive Open System ARchitecture}
\newacronym{bsw}{BSW}{Basic Software}
\newacronym{rte}{RTE}{Runtime Environment}
\newacronym{swc}{SWC}{Software Component}
\newacronym{ecu}{ECU}{Electronic Control Unit}
\newacronym{ota}{OTA}{Over-The-Air}
\newacronym{fusa}{FuSa}{Functional Safety}
\newacronym{asil}{ASIL}{Automotive Safety Integrity Level}
\newacronym{api}{API}{Application Programming Interface}
\newacronym{can}{CAN}{Controller Area Network}
\newacronym{lin}{LIN}{Local Interconnect Network}
\newacronym{som}{SoM}{System-on-Module}
\newacronym{os}{OS}{Operating System}
\newacronym{rtos}{RTOS}{Real-Time Operating System}
\newacronym{hal}{HAL}{Hardware Abstraction Layer}
\newacronym{llm}{LLM}{Large Language Model}
\newacronym{sdk}{SDK}{Software Development Kit}
\newacronym{ide}{IDE}{Integrated Development Environment}
\newacronym{ci}{CI}{Continuous Integration}
\newacronym{cd}{CD}{Continuous Deployment}
\newacronym{vm}{VM}{Virtuelle Maschine}
\newacronym{gc}{GC}{Garbage Collector}
\newacronym{abi}{ABI}{Application Binary Interface}
\newacronym{ffi}{FFI}{Foreign Function Interface}
\newacronym{iso}{ISO}{International Organization for Standardization}
\newacronym{iec}{IEC}{International Electrotechnical Commission}
\newacronym{misra}{MISRA}{Motor Industry Software Reliability Association}
\newacronym{mcal}{MCAL}{Microcontroller Abstraction Layer}
\newacronym{cp}{CP}{Classic Platform}
\newacronym{ap}{AP}{Adaptive Platform}
\newacronym{ara}{ARA}{AUTOSAR Runtime for Adaptive Applications}
\newacronym{posix}{POSIX}{Portable Operating System Interface}
\newacronym{adas}{ADAS}{Advanced Driver Assistance Systems}
\newacronym{v2x}{V2X}{Vehicle-to-Everything}
\newacronym{isr}{ISR}{Interrupt-Service-Routine}

\title{\textbf{Python und MicroPython als standardisierte Programmiersprache in AUTOSAR}}
\author{Stephan Epp}
\date{26. März 2026}

% ============================================================
\begin{document}
	
\maketitle
\tableofcontents
\clearpage

% ================================================================
\section{Einleitung}
\label{sec:einleitung}
Die Automobilindustrie steht an einem tiefgreifenden softwaretechnischen Wendepunkt.
Software-defined Vehicles, zunehmende Fahrzeugelektronik, \gls{adas}-Funktionen,
Over-the-Air-Updates und KI-basierte Regelkreise erfordern eine drastische
Steigerung der Softwareproduktivität sowie eine breitere Entwicklerbasis.
Innerhalb des \gls{autosar}-Konsortiums wird derzeit intensiv diskutiert,
die Programmiersprache Rust als offiziellen Bestandteil der Plattformspezifikationen
zu verankern~\cite{autosar_rust_wp2023, rust_automotive_safety2022}.

Diese Arbeit argumentiert auf Basis technischer, ökonomischer, ökosystemischer und
sicherheitstechnischer Analyse, dass \emph{Python} -- insbesondere in der Ausprägung
\emph{MicroPython}~\cite{micropython_main} -- die strategisch überlegene Wahl für
eine AUTOSAR-Standardisierung darstellt.
Python bietet die weltgrößte Entwicklercommunity, konkurrenzlose
Ökosystemreife, erstklassige KI/ML-Integration, ausgezeichnete Testbarkeit
und etablierte Einsatznachweise in sicherheitskritischen Systemen.
MicroPython ermöglicht dabei den direkten Einsatz auf ressourcenbeschränkten
Mikrocontrollern, wie sie in der \gls{autosar}-Classic-Plattform typisch sind,
während CPython in der Adaptive Platform nahtlos integrierbar ist.

Wir entwerfen eine vollständige Standardisierungsarchitektur, definieren
Sprachprofilstufen (PySafety Level 1--4), schlagen einen Migrationspfad vor
und diskutieren offene Forschungsfragen. Das Ergebnis ist ein konkreter,
umsetzbarer Fahrplan für die Aufnahme von Python/MicroPython als normierte
Sprache in AUTOSAR~R24 und darüber hinaus.

Die Geschichte der Automobilelektronik ist eine kontinuierliche Eskalation
der Softwarekomplexität. Während ein Fahrzeug der 1980er-Jahre kaum
Software enthielt, umfassen moderne Fahrzeuge über 100 \glspl{ecu},
mehr als 100~Millionen Zeilen Quellcode und ein Dutzend unterschiedliche
Kommunikationsbusse~\cite{pretschner2007software, broy2006challenges}.
Der Übergang zum Software-defined Vehicle (SDV) intensiviert diesen Trend
weiter: Fahrzeugfunktionen sollen künftig primär durch Software definiert,
aktualisiert und monetarisiert werden~\cite{sdv_mckinsey2023, sdv_bosch2024}.

Vor diesem Hintergrund gewinnt die Frage der Programmiersprachen-Standardisierung
innerhalb des \gls{autosar}-Rahmens strategische Bedeutung.
\gls{autosar} -- gegründet 2003 von BMW, Bosch, Continental, DaimlerChrysler,
Siemens VDO und Volkswagen -- definiert eine offene, standardisierte
Softwarearchitektur für Steuergeräte im Automobil~\cite{autosar_overview2017}.
Die Classic Platform (\gls{cp}) zielt auf Mikrocontroller mit begrenzten
Ressourcen; die 2017 eingeführte Adaptive Platform (\gls{ap}) adressiert
leistungsfähige Hochleistungsrechner mit POSIX-kompatiblem Betriebssystem~\cite{autosar_ap_explained2019}.

Beide Plattformen spezifizieren bislang C (für CP) bzw.\ C++14/17 (für AP)
als primäre Implementierungssprachen.
Angesichts bekannter Schwächen dieser Sprachen bezüglich
Speichersicherheit~\cite{memory_safety_nsa2022} formierte sich eine Gemeinschaft,
die Rust als modernere, speichersichere Alternative in AUTOSAR einzubringen
sucht~\cite{autosar_rust_wp2023, autosar_wg_rustsafety}.
Rust besitzt zweifellos Vorzüge in Bezug auf Speichersicherheit durch sein
Ownership-Modell~\cite{rust_book2022}, jedoch bringt es erhebliche Hindernisse mit sich:
steile Lernkurve, kleine Entwicklercommunity im Automotive-Bereich,
unreifes Tooling und fehlende Normnachweise für sicherheitskritische Anwendungen.

\textbf{Diese Arbeit positioniert Python/MicroPython als die überlegene Alternative.}
Während Rust eine nische Systemprogrammiersprache bleibt, ist Python die
meistgenutzte Programmiersprache der Welt~\cite{tiobe2025, stackoverflow_survey2024}
mit einem Ökosystem, das Machine Learning, formale Verifikation,
Testautomatisierung, Protokoll-Scripting und Hardware-nahe Programmierung
gleichermaßen abdeckt.
MicroPython~\cite{micropython_main, george2014micropython} -- eine schlanke,
auf Mikrocontroller optimierte Python-3-Implementierung -- beweist täglich
in Millionen industrieller und akademischer Anwendungen, dass Python
auch auf ressourcenbeschränkter Hardware praktikabel ist.

Die Forschungsfragen dieser Arbeit lauten:
\begin{enumerate}
  \item Welche technischen Anforderungen stellt \gls{autosar} an eine Programmiersprache,
        und inwiefern erfüllt Python/MicroPython diese?
  \item Wie lässt sich ein Python-Sprachprofil für AUTOSAR sicherheitstechnisch
        spezifizieren, das ISO~26262 konform ist?
  \item Welche Architektur ermöglicht die Integration von MicroPython in die
        AUTOSAR Classic Platform und Python in die Adaptive Platform?
  \item Welche wirtschaftlichen und ökosystemischen Vorteile entstehen durch
        eine Python-Standardisierung?
  \item Wie verhält sich Python/MicroPython im Vergleich zu Rust hinsichtlich
        der für AUTOSAR relevanten Kriterien?
\end{enumerate}

Der Rest dieser Arbeit ist wie folgt strukturiert:
Abschnitt~\ref{sec:autosar} erläutert die \gls{autosar}-Architektur im Detail.
Abschnitt~\ref{sec:python_micropython} stellt Python und MicroPython vor.
Abschnitt~\ref{sec:rust_analyse} analysiert die Rust-Bewegung kritisch.
Abschnitt~\ref{sec:technische_eignung} untersucht die technische Eignung.
Abschnitt~\ref{sec:sicherheit} behandelt Functional Safety und Cybersecurity.
Abschnitt~\ref{sec:architektur} entwirft die Integrationsarchitektur.
Abschnitt~\ref{sec:oekosystem} diskutiert das Ökosystem.
Abschnitt~\ref{sec:wirtschaft} analysiert den wirtschaftlichen Nutzen.
Abschnitt~\ref{sec:standardisierung} formuliert den Standardisierungsvorschlag.
Abschnitt~\ref{sec:verwandte} diskutiert verwandte Arbeiten.
Abschnitt~\ref{sec:fazit} schließt mit Fazit und Ausblick.

% ================================================================
\section{AUTOSAR -- Architektur und Anforderungen}
\label{sec:autosar}

\subsection{Geschichte und Motivation}
\gls{autosar} wurde 2003 als industrielles Konsortium gegründet~\cite{autosar_overview2017},
um die wachsende Fragmentierung der Steuergeräte-Software zu beenden.
Vor AUTOSAR pflegte jeder OEM und jeder Tier-1-Zulieferer eigene,
inkompatible Software-Stacks für jede Derivatefamilie.
Code war nicht portierbar, Integrationsprojekte hochkostenintensiv,
und Fehler von Zulieferern schwer zu isolieren.

AUTOSAR definierte eine layered architecture mit klar spezifizierten
Schnittstellen: Microcontroller Abstraction Layer (\gls{mcal}),
\gls{bsw}, \gls{rte} und Anwendungs-\glspl{swc}~\cite{autosar_layered2016}.
Durch standardisierte \glspl{api} und eine XML-basierte Beschreibungssprache
(ARXML) konnten Softwarekomponenten erstmals modular ausgetauscht werden.

\subsection{AUTOSAR Classic Platform (CP)}
\label{subsec:cp}

Die Classic Platform zielt auf Mikrocontroller der Klasse:
\begin{itemize}
  \item RAM: 32~kB bis 4~MB
  \item Flash: 256~kB bis 16~MB
  \item CPU: 8- bis 32-Bit, 8--800~MHz
\end{itemize}

CP spezifiziert C (ISO~C99, ergänzt durch MISRA~C~\cite{misra_c2012})
als Implementierungssprache. Der AUTOSAR~CP-Stack umfasst über 80~Module,
darunter OS (nach OSEK~\cite{osek_spec}), \gls{can}-Stack, Diagnostics (UDS),
Flash Management und mehr~\cite{autosar_cp_r23}.

\subsection{AUTOSAR Adaptive Platform (AP)}
\label{subsec:ap}

Die Adaptive Platform adressiert leistungsstarke Service-orientierte Architekturen
auf Hochleistungs-SoCs (z.B.\ NXP S32G, Renesas R-Car) mit:
\begin{itemize}
  \item Multi-Core ARM-A-Prozessoren, 1--16 Kerne, 1--3~GHz
  \item RAM: 1~GB bis 16~GB
  \item Ethernet/ SOME/IP als primärem Kommunikationsbus
  \item POSIX-Betriebssystem (z.B.\ Linux, QNX, Integrity)
\end{itemize}

AP nutzt C++14/17 als Sprache und definiert die
AUTOSAR Runtime for Adaptive Applications (\gls{ara})~\cite{autosar_ap_explained2019}.
Dienste (Service-oriented Communication via SOME/IP) ersetzen weitgehend
das Signal-orientierte Paradigma von CP.

\subsection{Anforderungskatalog an eine Programmiersprache in AUTOSAR}
\label{subsec:anforderungen}

Aus den AUTOSAR-Spezifikationen und den einschlägigen Normen
ISO~26262~\cite{iso26262_2018} und ISO~21434~\cite{iso21434_2021}
lassen sich folgende Anforderungen an eine Programmiersprache ableiten:

\begin{longtable}{lp{8cm}l}
\caption{Anforderungskatalog für Programmiersprachen in AUTOSAR}
\label{tab:anforderungen}\\
\toprule
\textbf{ID} & \textbf{Anforderung} & \textbf{Quelle} \\
\midrule
\endfirsthead
\multicolumn{3}{c}{\small\itshape (Fortsetzung von vorheriger Seite)}\\
\toprule
\textbf{ID} & \textbf{Anforderung} & \textbf{Quelle} \\
\midrule
\endhead
\bottomrule
\endfoot
R1  & Deterministisches Laufzeitverhalten; keine unvorhersehbaren Latenzspitzen & ISO~26262 \\
R2  & Qualifizierter Compiler / Interpreter (Tool Confidence Level) & ISO~26262-8 \\
R3  & Statische Analyse möglich (Typprüfung, Kontrollfluss) & MISRA~C \\
R4  & Definiertes Speichermodell; keine unkontrollierten Heap-Operationen & ISO~26262 \\
R5  & Sprachsubset definierbar für ASIL-B/C/D-Anwendungen & ISO~26262-6 \\
R6  & Formale oder halbformale Verifikation unterstützt & IEC~61508 \\
R7  & Breite Toolchain-Unterstützung (Debugger, Profiler) & Praxis \\
R8  & Hardware-nahe Programmierung möglich (Register, DMA, Interrupts) & CP-Anforderung \\
R9  & Interoperabilität mit C/C++ (FFI) & Migration \\
R10 & Ausreichend große, in Automotive erfahrene Entwicklercommunity & OEM-Anforderung \\
R11 & Offene Lizenzierung (keine proprietären Lock-ins) & AUTOSAR-Statut \\
R12 & Laufzeitfehler-Erkennung (Exceptions, Assertions) & ISO~26262-6 \\
R13 & Deterministische Speicherbereinigung kontrollierbar & Embedded \\
R14 & Cybersecurity-relevante Sprachkonstrukte (keine Buffer Overflows) & ISO~21434 \\
R15 & Testbarkeit (Unit, Integration, Mutationstest) & ISO~26262-6 \\
\end{longtable}

% ================================================================
\section{Python und MicroPython -- Überblick und Eigenschaften}
\label{sec:python_micropython}

\subsection{Python -- Die universelle Hochsprache}
\label{subsec:python}

Python wurde 1991 von Guido van Rossum veröffentlicht~\cite{vanrossum1995python}
und hat sich zur meistgenutzten Programmiersprache der Welt entwickelt~\cite{tiobe2025}.
Die Sprache vereint Prinzipien von Lesbarkeit, Ausdrucksstärke und
Produktivität. Charakteristisch sind:

\begin{itemize}
  \item \textbf{Dynamisches Typsystem} mit optionalem statischem Typing (seit PEP~484,
        Python~3.5) durch Typ-Annotationen und das \texttt{typing}-Modul~\cite{pep484};
  \item \textbf{Garbage Collected Memory Management} mit Referenzzählung als primärem
        Mechanismus (deterministischer als Mark-and-Sweep für kurzlebige Objekte);
  \item \textbf{Klare, erzwungene Syntaxstruktur} durch einrückungsbasierte Blöcke;
  \item \textbf{Erstklassige Funktionen, Closures, Generatoren} und
        \textbf{async/await}-Coroutinen (seit Python~3.5);
  \item \textbf{Umfangreiches Standardbibliotheks-Ökosystem} (über 200~000 Pakete
        auf PyPI~\cite{pypi_stats2024});
  \item \textbf{C-Erweiterbarkeit} via CPython C-API oder Cython~\cite{cython_book};
  \item \textbf{Vollständige Unicode-Unterstützung} (PEP~3120).
\end{itemize}

\subsubsection{Statische Typprüfung und Verifikation}
Ein weitverbreiteter Kritikpunkt an Python ist das dynamische Typsystem.
Dieser ist durch moderne Werkzeuge weitgehend obsolet: \texttt{mypy}~\cite{mypy_lehtosalo2016},
\texttt{pyright}~\cite{pyright_microsoft}, \texttt{pytype} (Google)
und \texttt{pyre} (Meta) bieten vollständige statische Typprüfung.
Kombiniert mit strikten Typ-Annotationen (PEP~484, PEP~526, PEP~544)
lässt sich Python-Code auf einem Niveau typprüfen, das C++
mit modernen Compilern vergleichbar ist~\cite{gros2021python_types}.

Für sicherheitskritische Anwendungen ist dies entscheidend:
Ein AUTOSAR-SWC könnte mit vollständiger Typ-Annotation versehen werden,
und die statische Analyse via mypy/pyright würde Typfehler bereits
vor dem Lauf entdecken, analog zu einem C-Compiler mit Warnungen.

\subsubsection{CPython und alternative Implementierungen}
Neben der Referenzimplementierung CPython existieren:
\begin{itemize}
  \item \textbf{MicroPython}: Für Mikrocontroller (s.\ Abschnitt~\ref{subsec:micropython});
  \item \textbf{PyPy}: JIT-kompilierter Interpreter, bis zu 5$\times$ schneller als CPython~\cite{pypy_bolz2009};
  \item \textbf{Nuitka}~\cite{nuitka}: Kompiliert Python zu nativem C-Code;
  \item \textbf{Cython}: Superset von Python, kompiliert zu C~\cite{cython_book};
  \item \textbf{mypyc}: Kompiliert typierten Python-Code zu C-Erweiterungen~\cite{mypyc2022}.
\end{itemize}

Insbesondere \textbf{Nuitka} und \textbf{mypyc} sind für die AUTOSAR-Adaptive-Platform
hochrelevant: Sie ermöglichen die Erstellung von nativen ELF-Binaries aus
Python-Quellcode, die direkt in POSIX-ECUs deploybar sind, ohne
einen Interpreter im Produktionssystem laufen zu lassen.

\subsection{MicroPython -- Python für Mikrocontroller}
\label{subsec:micropython}

MicroPython wurde 2013--2014 von Damien P. George entwickelt und
durch eine Kickstarter-Kampagne finanziert~\cite{george2014micropython}.
Es ist eine vollständige, schlanke Python-3-Implementierung
speziell für ressourcenbeschränkte Hardware~\cite{micropython_main}:

\begin{itemize}
  \item \textbf{Minimaler Flash-Bedarf}: ca.\ 256~kB (Basis-Interpreter);
        typische Deployment-Größe 256--512~kB Flash;
  \item \textbf{Minimaler RAM-Bedarf}: lauffähig ab 16~kB RAM,
        typisch 64--256~kB;
  \item \textbf{Python-3-Kompatibilität}: Kern-Sprachfeatures vollständig
        (Klassen, Generatoren, Closures, try/except, async/await);
  \item \textbf{Hardware-\glspl{api}}: \texttt{machine.Pin}, \texttt{machine.UART},
        \texttt{machine.SPI}, \texttt{machine.I2C}, \newline \texttt{machine.Timer},
        \texttt{machine.ADC}, \texttt{machine.WDT} u.v.m.;
  \item \textbf{Inline-Assembler} für ARM Thumb2, Xtensa und STM32;
  \item \textbf{Native Code Emitter}: \texttt{@micropython.native}-Decorator
        kompiliert zu native Maschinencode (4--10$\times$ schneller);
  \item \textbf{Viper-Code-Emitter}: \texttt{@micropython.viper}-Decorator
        mit C-ähnlichem Typsystem für höchste Performance;
  \item \textbf{Integrierter Garbage Collector} (Mark-and-Sweep mit
        konfigurierbarer Heap-Größe und manuell triggerbarer Kollektion).
\end{itemize}

\subsubsection{Unterstützte Mikrocontroller-Plattformen}
MicroPython unterstützt eine breite Plattformpalette~\cite{micropython_ports2023}:

\begin{table}[h]
\centering
\caption{Von MicroPython offiziell unterstützte Plattformen (Auswahl)}
\label{tab:micropython_platforms}
\begin{tabular}{llll}
\toprule
\textbf{Port} & \textbf{Architektur} & \textbf{Typische Automotive-MCU} & \textbf{Status} \\
\midrule
stm32  & ARM Cortex-M & STM32F4/F7/H7, Infineon TC3xx & Stabil \\
esp32  & Xtensa LX7, RISC-V & ESP32-S3, C6 & Stabil \\
rp2    & ARM Cortex-M0+ & RP2040, RP2350 & Stabil \\
nrf    & ARM Cortex-M & nRF52840 & Stabil \\
renesas-ra & ARM Cortex-M & RA4/RA6 & Beta \\
samd   & ARM Cortex-M0+ & SAMD51 & Stabil \\
unix   & x86/ARM64 & Linux-ECU (AP) & Stabil \\
windows & x86\_64 & Simulation/HIL & Stabil \\
\bottomrule
\end{tabular}
\end{table}

\subsubsection{MicroPython Performance}
Die Performance-Eigenschaften von MicroPython sind für eingebettete Systeme
gut untersucht~\cite{micropython_perf_papanastasiou2021, python_embedded_benchmark2022}:

\begin{itemize}
  \item Standardmäßiger bytecode-interpretierter Python-Code:
        ca.\ 10--50$\times$ langsamer als äquivalenter C-Code;
  \item Mit \texttt{@micropython.native}: ca.\ 2--5$\times$ langsamer als C;
  \item Mit \texttt{@micropython.viper}: annähernd C-Geschwindigkeit (1--2$\times$);
  \item Inline-Assembler: gleichwertig zu handoptimiertem Assembler.
\end{itemize}

Dies ist für nicht-zeitkritische Schichten von AUTOSAR-Software
(Diagnose, Konfiguration, OTA-Management, Kommunikationsprotokoll-Glue-Code)
vollkommen ausreichend und für zeitkritische Schichten durch
native/viper-Dekoratoren oder C-Extensions lösbar.

\subsubsection{MicroPython Speicherverwaltung im Detail}
\label{subsubsec:gc}

Das Speicherverwaltungsmodell von MicroPython ist für Embedded Systems ausgelegt:

\begin{itemize}
  \item \textbf{Heap-Allokation} nur aus einem konfigurierbaren, statisch reservierten
        Speicherbereich (kein malloc/free aus dem OS-Heap);
  \item \textbf{Deterministisches GC}: GC kann deaktiviert und manuell über
        \texttt{gc.collect()} ausgelöst werden;
  \item \textbf{Worst-Case Execution Time (WCET)}: Mit deaktiviertem GC und
        viper-Decorator ist WCET analysierbar;
  \item \textbf{Pre-Allokation}: Kritische Objekte können zur Boot-Zeit
        pre-allokiert und aus dem GC ausgeschlossen werden (\texttt{micropython.const});
  \item \textbf{Frozen Modules}: Python-Module können als Bytecode in Flash
        eingefroren werden (kein RAM für Quellcode nötig).
\end{itemize}

\begin{lstlisting}[language=MicroPython,caption={MicroPython Speichermanagement für zeitkritische Funktionen},label=lst:gc_control]
import gc
import micropython

# Pre-allokiere kritische Puffer zur Boot-Zeit
_can_rx_buffer = bytearray(64)
_can_tx_buffer = bytearray(64)

# Deaktiviere GC fuer zeitkritische Interrupt-Handler
@micropython.native
def can_isr_handler(frame_id: int, data: memoryview) -> None:
    """CAN-Interrupt-Handler: GC-frei, deterministisch."""
    gc.disable()  # Kein GC waehrend ISR
    try:
        _can_rx_buffer[:len(data)] = data
        _process_frame_native(frame_id, _can_rx_buffer)
    finally:
        gc.enable()

@micropython.viper
def _process_frame_native(frame_id: int, buf: ptr8) -> int:
    """Viper-Code: annaehernd C-Geschwindigkeit."""
    if frame_id == 0x100:
        return int(buf[0]) << 8 | int(buf[1])
    return -1
\end{lstlisting}

% ================================================================
\section{Analyse der Rust-Standardisierungsbewegung in AUTOSAR}
\label{sec:rust_analyse}

\subsection{Die Rust-Bewegung im Überblick}
Rust~\cite{rust_book2022} wurde von Mozilla Research entwickelt und 2015
in Version~1.0 veröffentlicht. Seine Haupteigenschaft ist das Ownership-System,
das Speichersicherheit ohne Garbage Collector zur Compilezeit garantiert.
Dies macht Rust attraktiv für Systemprogrammierung.

Im Automobilbereich wurde eine AUTOSAR-Arbeitsgruppe für Rust gebildet~\cite{autosar_wg_rustsafety},
und das Whitepaper ``Safe Rust Guidelines for Automotive~\cite{autosar_rust_wp2023}''
beschreibt erste Richtlinien. Die MISRA-Community diskutiert MISRA~Rust~\cite{misra_rust2022}.

\subsection{Kritische Analyse der Rust-Eignung für AUTOSAR}
\label{subsec:rust_kritik}

Trotz dieser Aktivitäten bestehen fundamentale Hindernisse:

\subsubsection{Lernkurve und Entwicklerengpass}
Rust gilt als eine der schwierigsten Programmiersprachen~\cite{stackoverflow_survey2024}.
Das Ownership-/Borrow-Modell ist konzeptuell neuartig und erfordert
erhebliche Einarbeitungszeit.
Während Python nach TIOBE~\cite{tiobe2025} und dem
Stack Overflow Survey~\cite{stackoverflow_survey2024} weltweit die
meistgenutzte Sprache ist (ca.\ 30~Mio. aktive Entwickler), zählt
Rust nur ca.\ 3~Mio. Entwickler weltweit -- davon ein Bruchteil
mit Automotive-Erfahrung. Der resultierende Entwicklermangel
verteuert Projekte und verlängert Time-to-Market erheblich.

\subsubsection{Fehlendes Safety-Zertifizierungs-Ökosystem}
ISO~26262-8 verlangt für Sprachtools einen
\emph{Tool Confidence Level}~\cite{iso26262_2018}.
Für Rust existiert (Stand 2025) kein zertifizierter Compiler
(nach ISO~26262 ASIL-D qualifiziert), während GCC und LLVM/Clang
seit Jahren entsprechend qualifiziert werden (z.B.\ durch
AbsInt, AdaCore, ETAS).
MicroPython und CPython können auf qualifizierte C-Backends
(LLVM/Clang, GCC mit Qualifizierungskit) aufsetzen.

\subsubsection{Fehlende AUTOSAR-Treiberintegration}
Automotive-spezifische HAL-Bibliotheken (CAN, LIN, FlexRay, Ethernet,
diagnostischer Speicherzugriff) existieren für C über Jahrzehnte.
Für Rust sind automotive-grade BSW-Bindungen weitgehend nicht existent
und müssten von Grund auf neu entwickelt werden -- ein
Jahrzehnt-Projekt für die gesamte Industrie.

\subsubsection{`unsafe` Rust: kein Gewinn gegenüber C}
Viele Hardware-nahe Operationen in Automotive-ECUs erfordern
zwingend \texttt{unsafe}-Rust-Blöcke~\cite{rust_unsafe2021}.
Damit sind die Speichersicherheits-Garantien von Rust in den
kritischsten Schichten (MCAL, Treiber) ohnehin nicht gegeben.
Python/MicroPython mit seiner klar definierten Sicherheitsschicht
(C-Extension in sicherem C) ist hier transparent.

\subsubsection{Zusammenfassung: Rust vs. Python/MicroPython}

\begin{table}[h]
\centering
\caption{Vergleich Rust vs.\ Python/MicroPython für AUTOSAR-Standardisierung}
\label{tab:rust_vs_python}
\begin{tabularx}{\textwidth}{lXX}
\toprule
\textbf{Kriterium} & \textbf{Rust} & \textbf{Python/MicroPython} \\
\midrule
Entwicklerbasis weltweit & $\sim$3~Mio. & $>$30~Mio. \\
Automotive-Erfahrung (Community) & sehr gering & mittel--hoch (Skripte, Tests, Tools) \\
Lernkurve & sehr steil & flach \\
Compiler-Qualifizierung (ISO~26262) & nicht vorhanden & via GCC/LLVM möglich \\
AUTOSAR BSW-Bindungen & keine & C-FFI, erste Bibliotheken \\
MISRA-äquivalentes Profil & MISRA Rust (Draft) & PyMISRA-Ansätze (forschbar) \\
Embedded-Footprint & 50--500~kB (no\_std) & 256~kB (MicroPython) \\
KI/ML-Integration & rudimentär & erstklassig (PyTorch, TensorFlow) \\
OTA-Scripting & nicht vorgesehen & nativ via Python-Scripts \\
Testökosystem & wachsend & ausgereift (pytest, hypothesis) \\
Statische Analyse & sehr gut & sehr gut (mypy, pyright) \\
Formale Verifikation & experimentell & SPARK-ähnlich via typsichere Teilmengen \\
Hardware-nahe Ops & native (unsafe) & viper-Decorator, Inline-ASM \\
Industrielle Nachweise & kaum & MicroPython in Industrie weit verbreitet \\
\bottomrule
\end{tabularx}
\end{table}

% ================================================================
\section{Technische Eignung von Python/MicroPython für AUTOSAR}
\label{sec:technische_eignung}

\subsection{Erfüllung des Anforderungskatalogs}
\label{subsec:anforderungen_erfuellung}

Wir prüfen systematisch, inwieweit Python/MicroPython die in
Tabelle~\ref{tab:anforderungen} definierten Anforderungen erfüllt:

\paragraph{R1 -- Deterministisches Laufzeitverhalten}
MicroPython erlaubt die Kontrolle des GC, die Verwendung von
\texttt{@micropython.native} und \texttt{@micropython.viper} sowie
Inline-Assembler für zeitkritische Pfade.
Durch Pre-Allokation kritischer Objekte und manuelle GC-Steuerung
(vgl.\ Listing~\ref{lst:gc_control}) ist deterministisches Verhalten
erzielbar. Für ASIL-D-Funktionen wäre eine static-alloc-only-Bibliothek
analog zu \texttt{heapless} in Rust zu spezifizieren.

\paragraph{R2 -- Qualifizierter Interpreter}
CPython kompiliert zu C-Code, der mit qualifizierten GCC/Clang-Compilern
übersetzt wird. MicroPython wird ebenfalls mit GCC/Clang kompiliert.
Eine Qualifizierung des MicroPython-Interpreters selbst nach
Tool Confidence Level (TCL2) ist machbar -- vergleichbar mit
der Qualifizierung von Java~VM-Implementierungen nach IEC~61508~\cite{jvm_qualification2018}.

\paragraph{R3 -- Statische Analyse}
mypy mit strikter Konfiguration (\texttt{--strict}) und pyright
ermöglichen vollständige statische Typprüfung. Zusätzlich bieten
\texttt{pylint}~\cite{pylint2023}, \texttt{flake8}, \texttt{bandit}
(Sicherheitsanalyse) und \texttt{pyflakes} umfangreiche Codeanalyse.
Für AUTOSAR-spezifische Regeln wäre ein \emph{AUTOSAR Python Checker}
als Open-Source-Tool definierbar.

\paragraph{R4 -- Definiertes Speichermodell}
MicroPython: gesamter Heap ist ein konfigurierbarer, statischer Buffer.
Keine Interaktion mit OS-Heap. Vollständig kontrollierbar.
Für Adaptive Platform: Python mit speziell allokierten arenas
(CPython configurierbar via \texttt{pymalloc} oder
custom allocator~\cite{cpython_allocator2022}).

\paragraph{R5 -- Sprachsubset für ASIL}
Python besitzt eine klar definierte, kleine Kernsprache.
Ein ASIL-Sprachsubset (``PySafety-Profile'') ist eindeutig definierbar:
Verboten wären z.B.\ \texttt{eval()}, \texttt{exec()}, dynamischer Import,
und unbeschränkte Rekursion. Erlaubt wären typisierte Funktionen,
Pre-allokierte Datenstrukturen und native-Dekoratoren.
(s.~Abschnitt~\ref{sec:sprachprofil})

\paragraph{R6 -- Formale Verifikation}
Für typisierte Python-Teilmengen existieren Ansätze zu
Model Checking (via Abstraktion auf Typsystem-Ebene)~\cite{gros2021python_types}.
Das Projekt \emph{nagini}~\cite{nagini2018} implementiert
Verifikation von Python-Code via Viper-Assertion-Sprache.

\paragraph{R7 -- Toolchain}
Python verfügt über ausgezeichnete Toolchain-Unterstützung:
Eclipse/VS~Code-Integration, REPL für interaktives Debugging,
GDB-Integration via \texttt{python-\newline gdb.py}~\cite{cpython_gdb},
Profiler (\texttt{cProfile, line\_profiler}), Memory-Analyzer (\texttt{tracemalloc}).

\paragraph{R8 -- Hardware-nahe Programmierung}
MicroPython bietet native Hardware-APIs. Listing~\ref{lst:hw_access}
zeigt exemplarisch den direkten Registerzugriff auf einem Cortex-M-Controller.

\begin{lstlisting}[language=MicroPython,caption={Direkter Hardware-Registerzugriff mit MicroPython},label=lst:hw_access]
import micropython
from micropython import const

# Direkte Speicher-mapped Register-Adresse (Cortex-M)
_RCC_BASE  = const(0x40023800)
_RCC_AHB1ENR = const(_RCC_BASE + 0x30)
_GPIOA_MODER = const(0x40020000)

@micropython.viper
def enable_gpioa_clock() -> None:
    """Aktiviere GPIOA-Takt ueber RCC-Register (direkte Schreiboperation)."""
    # Zeiger auf uint32-Register
    rcc_ahb1enr = ptr32(_RCC_AHB1ENR)
    rcc_ahb1enr[0] = rcc_ahb1enr[0] | 0x1  # Bit 0: GPIOAEN

@micropython.viper  
def config_pin_output(pin: int) -> None:
    """Konfiguriere GPIO-Pin als Push-Pull-Ausgang."""
    moder = ptr32(_GPIOA_MODER)
    moder[0] = (moder[0] & ~(0x3 << (pin * 2))) | (0x1 << (pin * 2))
\end{lstlisting}

\paragraph{R9 -- Interoperabilität mit C/C++}
MicroPython unterstützt native C-Module (\texttt{.c}-Dateien, die als
MicroPython-Module kompiliert werden)~\cite{micropython_cmodules2023}.
CPython bietet die vollständige C-API sowie \texttt{ctypes},
\texttt{cffi} und Cython. Bestehende AUTOSAR-BSW-C-Bibliotheken
können direkt von Python aus aufgerufen werden.

\paragraph{R10--R15}
Alle weiteren Anforderungen werden in den folgenden Abschnitten
(Sicherheit, Architektur, Ökosystem) detailliert behandelt.

\subsection{MicroPython auf Automotive-Mikrocontrollern}
\label{subsec:automotive_mcu}

Automotive-Mikrocontroller der Generation 2020+ haben typische Spezifikationen,
die MicroPython komfortabel unterstützen:

\begin{table}[h]
\centering
\caption{Ressourcenvergleich: MicroPython-Anforderungen vs.\ Automotive-MCUs}
\label{tab:mcu_resources}
\begin{tabular}{llll}
\toprule
\textbf{Ressource} & \textbf{MicroPython Min.} & \textbf{Renesas RH850/F1x} & \textbf{Infineon TC39x} \\
\midrule
Flash             & 256~kB           & 4~MB             & 16~MB \\
RAM               & 64~kB            & 512~kB           & 8~MB \\
CPU-Takt          & 8~MHz            & 240~MHz          & 300~MHz \\
CPU-Architektur   & ARM/RISC-V       & V850 / ARM       & TriCore \\
\bottomrule
\end{tabular}
\end{table}

Die Ressourcen moderner Automotive-MCUs übertreffen die MicroPython-Mindestanfor\-derungen
um einen Faktor von 16--64$\times$ (Flash) bzw.\ 8--128$\times$ (RAM).
Selbst ohne die gesamte Python-Stdlib ist genügend Platz für
eine vollwertige MicroPython-Laufzeit plus applikationsspezifischen Code.

\subsection{Latenz und Interrupt-Handling}

\label{subsec:latenz}

Interrupt-Service-Routinen (\glspl{isr}) sind der hardwarenächste Teil des Systems.
MicroPython bietet mehrere Möglichkeiten:

\begin{enumerate}
  \item \textbf{C-ISR}: Interrupt-Handler in C implementiert, ruft Python-Callback
        via \texttt{micropython.\newline schedule()} aus sicherem Kontext auf;
  \item \textbf{Native ISR}: \texttt{@micropython.native}-dekorierter Python-Handler,
        direkt als ISR registrierbar (\texttt{machine.Pin.irq(..., handler=fn)});
  \item \textbf{Deferred Execution}: ISR setzt Flag, Hauptschleife verarbeitet
        asynchron (klassisches Muster in embedded systems).
\end{enumerate}

Messungen auf STM32F4 (168~MHz Cortex-M4) zeigen~\cite{micropython_irq_bench2022}:
\begin{itemize}
  \item C-ISR Latenz: $<$1~µs
  \item Native Python-ISR Latenz: 2--5~µs
  \item Deferred Python-Callback: 10--50~µs (abhängig von Scheduling)
\end{itemize}

Für AUTOSAR-CP-Anwendungen mit typischen Zykluszeiten von
1~ms (Task~Class~1) bis 100~ms (Task~Class~3) sind diese Latenzen
vollkommen akzeptabel. Nur harte Echtzeit-ISRs ($<$1~µs) würden
weiterhin native C erfordern.

% ================================================================
\section{Sicherheitstechnische Betrachtung}
\label{sec:sicherheit}

\subsection{Functional Safety nach ISO~26262}
\label{subsec:fusa}

ISO~26262~\cite{iso26262_2018} definiert vier Automotive Safety Integrity Levels
(ASIL~A bis ASIL~D, aufsteigend kritisch). Die Auswahl der Programmiersprache
und ihrer Laufzeitumgebung fließt in den Safety Case ein.

\subsubsection{Anforderungen an die Sprachlaufzeit}
ISO~26262-6 (Software-Level) fordert:
\begin{itemize}
  \item Verwendung sprach-spezifischer Sprachguidelines
        (Subset, Coding Guidelines);
  \item Nachweis der Werkzeugqualifizierung;
  \item Statische Analyse, Kontrollflussanalyse, Datenflussanalyse;
  \item Dynamisches Testen (Statement Coverage ASIL-A/B,
        Branch Coverage ASIL-C, MC/DC ASIL-D).
\end{itemize}

\subsubsection{PySafety-Sprachprofil}
\label{sec:sprachprofil}

Wir definieren vier \textbf{PySafety-Level} als AUTOSAR-Erweiterung:

\begin{definition}[PySafety Level]
  Ein PySafety Level beschreibt ein standardisiertes Sprachsubset
  von Python/MicroPython, das für eine bestimmte ASIL-Klasse einsetzbar ist.
\end{definition}

\begin{longtable}{lp{5cm}p{5cm}l}
\caption{PySafety-Level-Definitionen}
\label{tab:pysafety}\\
\toprule
\textbf{Level} & \textbf{Erlaubte Sprachfeatures} & \textbf{Verbotene Features} & \textbf{ASIL-Mapping} \\
\midrule
\endfirsthead
\toprule
\textbf{Level} & \textbf{Erlaubte Features} & \textbf{Verbotene Features} & \textbf{ASIL} \\
\midrule
\endhead
\bottomrule
\endfoot
PSL-0 & Vollständiges Python & -- & QM \\
PSL-1 & Python ohne eval/exec, ohne dynamischen Import & eval, exec, \_\_import\_\_ & ASIL-A \\
PSL-2 & PSL-1 + vollständige Typ-Annotationen, kein \texttt{*args}/\texttt{**kwargs} ohne Typen & Dynamische Attribute, getattr ohne Typcheck & ASIL-B \\
PSL-3 & PSL-2 + nur pre-allokierte Container, kein unbeschränkter Heap & list(), dict() ohne Größenbeschränkung, offene Rekursion & ASIL-C \\
PSL-4 & PSL-3 + nur viper/native-Dekoratoren in ISR-Pfaden, belegter WCET-Nachweis & Alle GC-triggernden Operationen in zeitkritischen Pfaden & ASIL-D \\
\end{longtable}

\subsubsection{Coding Guidelines und statische Prüfung}
Analog zu MISRA~C~\cite{misra_c2012} wären \textbf{AUTOSAR~Python~Guidelines} zu definieren.
Erste Entwürfe könnten auf bestehenden Arbeiten aufbauen:
\begin{itemize}
  \item \textbf{PySafety Coding Rules}: Ein Satz von ca.\ 80--120 Regeln,
        die mit \texttt{pylint}-Plugin oder eigenem Checker automatisiert
        prüfbar sind;
  \item \textbf{NASA PySafe}: NASA's Python-Codierungsrichtlinien
        für sicherheitskritische Systeme~\cite{nasa_python2011};
  \item \textbf{CERT~Secure Coding Standard für Python}~\cite{cert_python2019}.
\end{itemize}

\subsection{Cybersecurity nach ISO~21434}
\label{subsec:cybersecurity}

ISO~21434~\cite{iso21434_2021} regelt Cybersecurity im Fahrzeugbereich.
Python bietet hier signifikante Vorteile gegenüber C/C++:

\begin{itemize}
  \item \textbf{Keine Buffer Overflows}: Python's Memory-Model verhindert
        klassische Stack- und Heap-Overflows on language level;
  \item \textbf{Keine Format-String-Vulnerabilities}: f-Strings und
        \texttt{str.format()} sind typsicher;
  \item \textbf{Integer-Overflow frei}: Python's \texttt{int} ist beliebig groß
        (in MicroPython: 30-bit integer intern, erweiterbar);
  \item \textbf{Kryptographie-Standardbibliothek}: \texttt{hashlib}, \texttt{hmac},
        \texttt{ssl} (CPython) oder \texttt{ucry\newline ptolib} (MicroPython);
  \item \textbf{Security Linter}: \texttt{bandit}~\cite{bandit2023} scannt
        Python-Code auf OWASP-Top-10-Schwachstellen.
\end{itemize}

Für Automotive gelten besondere Anforderungen:
\begin{itemize}
  \item UDS-Sicherheits-Services (Session-Management, Seed-Key-Algorithmus)
        in Python implementierbar mit kryptografischen Primitiven;
  \item SecureBoot-Verifikation via \texttt{ucryptolib.aes} auf MicroPython;
  \item TLS~1.3-Kommunikation via \texttt{ussl} auf MicroPython oder
        mbedTLS-Binding.
\end{itemize}

\begin{lstlisting}[language=MicroPython,caption={Sicherheits-relevante Kryptographie in MicroPython},label=lst:crypto]
import ucryptolib
import hashlib
import ubinascii

# AES-128-CBC fuer SecureFlash-Verifikation
def verify_firmware_signature(
    firmware: bytes,
    expected_hmac: bytes,
    key: bytes
) -> bool:
    """Pruefe HMAC-SHA256 einer Firmware (MicroPython)."""
    import hmac
    computed = hmac.new(key, firmware, hashlib.sha256).digest()
    # Constant-time Vergleich (timing-attack-safe)
    return _constant_time_compare(computed, expected_hmac)

def _constant_time_compare(a: bytes, b: bytes) -> bool:
    """Zeitkonstanter Byte-String-Vergleich."""
    if len(a) != len(b):
        return False
    result = 0
    for x, y in zip(a, b):
        result |= x ^ y
    return result == 0
\end{lstlisting}

\subsection{WCET-Analyse und Echtzeitfähigkeit}
\label{subsec:wcet}

Worst-Case Execution Time Analysis (WCET) ist für ASIL-C/D-Funktionen
zwingend erforderlich~\cite{wilhelm2008wcet}.
Für Python/MicroPython-Code ist WCET auf drei Ebenen möglich:

\begin{enumerate}
  \item \textbf{Viper-Level}: Code mit \texttt{@micropython.viper} ist
        strukturell äquivalent zu C-Code und kann mit etablierten
        WCET-Werkzeugen (AbsInt aiT~\cite{absint_ait}, TriCore WCET-Analyzer)
        analysiert werden -- der Viper-Code wird zu identischem Maschinencode
        wie äquivalentes C kompiliert;
  \item \textbf{Bytecode-Level}: Für bytecode-interpretierten Python-Code
        ist WCET durch Analyse des Interpreters + Bytecode berechenbar
        (Forschungsansatz: ~\cite{micropython_wcet2023});
  \item \textbf{Partitionierung}: Zeitkritische Funktionen in C, zeitunkritische
        in Python -- bewährter Praxis-Ansatz.
\end{enumerate}

% ================================================================
\section{Integrationsarchitektur}
\label{sec:architektur}

\subsection{Überblick: Python in der AUTOSAR-Schichtenarchitektur}
\label{subsec:schichten}

\begin{figure}[h]
\centering
% ASCII-Art Darstellung der Architektur (da kein externes Bild verfuegbar)
\begin{tabular}{|c|}
\hline
\rowcolor{autosarblue!30}
\textbf{AUTOSAR Application Layer (SWCs)} \\
\hline
\rowcolor{pythonblue!20}
\textbf{Python / MicroPython SWC Layer} \\
(Python-SWCs, typisiert, PSL-1..4) \\
\hline
\rowcolor{pythonblue!10}
\textbf{PySafety Runtime (PyAutosar-RTE)} \\
(Python-RTE-Schnittstelle, C-FFI-Bridge) \\
\hline
\rowcolor{gray!20}
\textbf{AUTOSAR RTE (Standard C)} \\
\hline
\rowcolor{gray!10}
\textbf{AUTOSAR BSW (Standard C)} \\
(COM, NM, DEM, DCM, NvM, ...) \\
\hline
\rowcolor{gray!5}
\textbf{MCAL (C, Hersteller-spezifisch)} \\
\hline
\rowcolor{lightgray}
\textbf{Hardware (ECU, MCU, Peripherie)} \\
\hline
\end{tabular}
\caption{Integrierte AUTOSAR-Schichtenarchitektur mit Python/MicroPython SWC-Ebene}
\label{fig:architektur}
\end{figure}

\subsection{Python SWC Integration -- Classic Platform}
\label{subsec:cp_integration}

Für die AUTOSAR Classic Platform wird folgende Integrationsarchitektur vorgeschlagen:

\subsubsection{PyAutosar-RTE-Bridge}
Eine C-Schicht (\texttt{pyautosar\_rte.c}) bildet die AUTOSAR-RTE-API
auf Python-callable C-Funktionen ab:

\begin{lstlisting}[language=C,caption={PyAutosar-RTE-Bridge: C-Seite (Auszug)},label=lst:rte_bridge_c]
/* pyautosar_rte.c: Bridge zwischen AUTOSAR RTE und MicroPython VM */
#include "Rte_Type.h"
#include "Rte_MySwc.h"
#include "py/runtime.h"  /* MicroPython C-API */

/* Python-callable: Rte_Write_Port_Signal */
STATIC mp_obj_t mp_Rte_Write_VehicleSpeed(mp_obj_t speed_obj) {
    float speed = mp_obj_get_float(speed_obj);
    Std_ReturnType ret = Rte_Write_MySwc_VehicleSpeed(
        (Float32)speed
    );
    return mp_obj_new_int(ret);
}
STATIC MP_DEFINE_CONST_FUN_OBJ_1(
    mp_Rte_Write_VehicleSpeed_obj,
    mp_Rte_Write_VehicleSpeed
);

/* Modul-Registrierung */
STATIC const mp_rom_map_elem_t autosar_rte_module_globals_table[] = {
    { MP_ROM_QSTR(MP_QSTR___name__), MP_ROM_QSTR(MP_QSTR_autosar_rte) },
    { MP_ROM_QSTR(MP_QSTR_write_vehicle_speed),
      MP_ROM_PTR(&mp_Rte_Write_VehicleSpeed_obj) },
};
\end{lstlisting}

\begin{lstlisting}[language=MicroPython,caption={Python SWC unter AUTOSAR CP},label=lst:python_swc]
"""
Beispiel: Python-implementierter AUTOSAR SWC fuer Fahrzeuggeschwindigkeit
PSL-3 compliant: Typ-Annotationen, keine dynamischen Allokierungen
"""
from autosar_rte import write_vehicle_speed, read_wheel_speeds
import micropython

# Kalibrierungsdaten: pre-allokiert, const
WHEEL_CIRC_MM: float = micropython.const(1975.0)  # 215/55R17
PULSES_PER_REV: int  = micropython.const(48)

@micropython.native
def runnable_calc_speed() -> None:
    """
    AUTOSAR Runnable: 10ms zyklisch aufgerufen.
    Liest Raddrehzahl-Pulse, berechnet Fahrzeuggeschwindigkeit.
    """
    pulse_fl: int
    pulse_fr: int
    pulse_rl: int
    pulse_rr: int
    pulse_fl, pulse_fr, pulse_rl, pulse_rr = read_wheel_speeds()

    # Durchschnitt Hinterachse (Antriebsachse)
    avg_pulses: int = (pulse_rl + pulse_rr) >> 1

    # Umrechnung: Pulse/10ms -> km/h
    # v [km/h] = (pulses * circ_mm) / (pulses_per_rev * 10ms) * 3.6
    speed_kmh: float = (
        avg_pulses * WHEEL_CIRC_MM * 3.6
    ) / (PULSES_PER_REV * 10.0)

    write_vehicle_speed(speed_kmh)
\end{lstlisting}

\subsubsection{MicroPython-VM Integration in AUTOSAR-OS}

Die MicroPython-\gls{vm} wird als eigenständiger AUTOSAR-Task unter
dem AUTOSAR-OS (OSEK-basiert) gestartet:

\begin{lstlisting}[language=C,caption={MicroPython VM als AUTOSAR-OS-Task},label=lst:mpy_task]
/* autosar_python_task.c */
#include "Os.h"
#include "py/compile.h"
#include "py/runtime.h"
#include "py/gc.h"

/* Statisch allokierter Heap fuer MicroPython VM */
static uint8_t mp_heap[32768];  /* 32 kB Python-Heap */

TASK(PythonRunnableTask) {
    /* Initialisierung (einmalig beim ersten Aufruf) */
    static bool initialized = false;
    if (!initialized) {
        gc_init(mp_heap, mp_heap + sizeof(mp_heap));
        mp_init();
        initialized = true;
        /* Lade Python-SWC aus gefrorenem Modul */
        mp_import_mod(MP_QSTR_swc_speed_calc);
    }
    /* Rufe Python-Runnable auf */
    mp_call_function_0(
        mp_load_global(MP_QSTR_runnable_calc_speed)
    );
    TerminateTask();
}
\end{lstlisting}

\subsection{Python Integration -- Adaptive Platform}
\label{subsec:ap_integration}

Auf der Adaptive Platform ist die Integration trivialer:
CPython (oder Nuitka-kompiliertes Python) läuft als nativer POSIX-Prozess.
Die ara-Schnittstellen (\gls{ara}) werden als Python-Bindings bereitgestellt:

\begin{lstlisting}[language=Python,caption={Python Adaptive Application (ara-Bindings)},label=lst:pyara]
"""
AUTOSAR Adaptive Platform: Python Service Application.
Kommuniziert via SOME/IP ueber ara::com Python-Bindings.
"""
import asyncio
from ara.com import ServiceProxy, ServiceSkeleton
from ara.log import Logger
from ara.exec import ApplicationClient
from typing import Optional

logger = Logger.create_logger("SpeedService")

class VehicleSpeedSkeleton(ServiceSkeleton):
    """AUTOSAR AP Service-Skeleton in Python."""
    
    async def get_speed(self) -> float:
        """Gibt aktuelle Fahrzeuggeschwindigkeit zurueck (km/h)."""
        return self._current_speed
    
    async def on_speed_update(self, speed: float) -> None:
        """Event-Handler: Neue Geschwindigkeit empfangen."""
        logger.info(f"Geschwindigkeitsupdate: {speed:.1f} km/h")
        self._current_speed = speed
        await self.notify_speed_changed(speed)

async def main() -> None:
    """AUTOSAR AP Application Main-Coroutine."""
    app = ApplicationClient("SpeedService")
    skeleton = VehicleSpeedSkeleton(
        instance_id="SpeedService/1"
    )
    
    async with app.run():
        await skeleton.offer_service()
        logger.info("SpeedService angeboten -- warte auf Clients")
        await asyncio.sleep_until_cancel()

if __name__ == "__main__":
    asyncio.run(main())
\end{lstlisting}

\subsection{pyautosar -- Existierende Python/AUTOSAR-Tooling-Infra-\\struktur}
\label{subsec:pyautosar}

Die Behauptung, Python sei fremd für das AUTOSAR-Ökosystem, ist falsch.
Das Open-Source-Projekt \texttt{autosar}~\cite{pyautosar_gh} (früher \texttt{pyautosar})
bietet bereits heute:

\begin{itemize}
  \item ARXML-Parsing und -Generierung in Python;
  \item SWC-Beschreibung via Python-API;
  \item Port-Interface-Generierung;
  \item Integration mit AUTOSAR-Buildtools.
\end{itemize}

Dieses Projekt belegt, dass die AUTOSAR-Community Python aktiv nutzt
und die Sprachintegration praktisch erprobt ist.

\subsection{OTA-Update via Python}
\label{subsec:ota}

Eines der wichtigsten Anwendungsszenarien für Python in AUTOSAR ist
das Over-the-Air-Update-Management~\cite{ota_autosar_ap2021}.
Python's Stärken kommen hier besonders zum Tragen:

\begin{lstlisting}[language=Python,caption={OTA-Update-Manager in Python (AUTOSAR AP)},label=lst:ota]
"""OTA Update Manager fuer AUTOSAR Adaptive Platform."""
import hashlib
import asyncio
from pathlib import Path
from ara.ucm import UpdateCampaignManagement, SoftwarePackage

class PythonOTAManager:
    """Verwaltet OTA-Updates gem. AUTOSAR UCM-Spezifikation."""
    
    def __init__(self, ucm: UpdateCampaignManagement) -> None:
        self._ucm = ucm
        self._download_dir = Path("/var/autosar/ota")
    
    async def handle_update_campaign(
        self,
        package: SoftwarePackage
    ) -> bool:
        """Fuehrt einen vollstaendigen OTA-Update-Zyklus durch."""
        # 1. Integritaetspruefung
        if not await self._verify_package(package):
            return False
        # 2. Platzbedarf pruefen
        if not await self._check_disk_space(package):
            return False
        # 3. Herunterladen + Verifizieren
        local_path = await self._download_package(package)
        # 4. UCM-Transfer
        return await self._ucm.install(local_path)
    
    async def _verify_package(
        self,
        pkg: SoftwarePackage
    ) -> bool:
        data = await pkg.get_manifest()
        digest = hashlib.sha256(data).hexdigest()
        return digest == pkg.expected_sha256
\end{lstlisting}

% ================================================================
\section{Das Python-Ökosystem im Automotive-Kontext}
\label{sec:oekosystem}

\subsection{Weltgrößtes Entwickler-Ökosystem}
\label{subsec:entwickler}

Python ist gemäß Stack Overflow Developer Survey 2024~\cite{stackoverflow_survey2024}
die meistgenutzte Programmiersprache unter professionellen
Entwicklern (ca.\ 51\% General-Use).
Der Python Package Index (PyPI)~\cite{pypi_stats2024} zählt über
500~000~Pakete (Stand 2025), darunter:

\begin{itemize}
  \item \textbf{KI/ML}: PyTorch, TensorFlow, scikit-learn, ONNX Runtime;
  \item \textbf{Datenwissenschaft}: NumPy, pandas, SciPy, matplotlib;
  \item \textbf{Netzwerk}: Twisted, asyncio, aiohttp, cryptography;
  \item \textbf{Testen}: pytest, hypothesis (Property-based Testing), unittest.mock;
  \item \textbf{Formale Methoden}: nagini, z3-python-bindings;
  \item \textbf{Automotive-spezifisch}: python-can~\cite{python_can2023},
        cantools, python-udsoncan, pyscapy (Fahrzeugdiagnose),
        autosar (ARXML-Tooling~\cite{pyautosar_gh}).
\end{itemize}

\subsection{Testautomatisierung}
\label{subsec:testing}

Python ist die Standardsprache für Testautomatisierung im Automotive-Bereich:

\begin{itemize}
  \item \textbf{pytest}~\cite{pytest2023}: Defacto-Standard für Unit- und
        Integrationstests; erweiterbar über Plugins;
  \item \textbf{hypothesis}~\cite{hypothesis2019}: Property-based Testing,
        äquivalent zu QuickCheck in Haskell;
  \item \textbf{Robot Framework}: Akzeptanztestrahmen, weit verbreitet in
        HIL-Testsystemen (Vector CANoe/CANalyzer-Anbindung);
  \item \textbf{tox}: Multi-Environment-Testmatrix;
  \item \textbf{coverage.py}: Measurement Coverage (Statement, Branch, MC/DC via
        Erweiterung).
\end{itemize}

Für AUTOSAR-Anwendungen bedeutet dies: Eine in Python implementierte
SWC kann mit denselben Testetools getestet werden wie das übrige
Python-Produktionscode-Ökosystem -- ohne zusätzliche Adapterschichten.

\begin{lstlisting}[language=Python,caption={pytest-basierter Unit-Test eines Python-AUTOSAR-SWC},label=lst:pytest]
"""Unit-Tests fuer den SpeedCalculation-SWC (pytest)."""
import pytest
from unittest.mock import MagicMock, patch
from swc_speed_calc import runnable_calc_speed, WHEEL_CIRC_MM, PULSES_PER_REV

@pytest.fixture
def mock_rte():
    """Mock der AUTOSAR-RTE-Schnittstelle."""
    with patch("swc_speed_calc.read_wheel_speeds") as mock_read, \
         patch("swc_speed_calc.write_vehicle_speed") as mock_write:
        yield mock_read, mock_write

class TestSpeedCalculation:
    def test_stillstand(self, mock_rte):
        """Fahrzeug steht: Geschwindigkeit = 0."""
        mock_read, mock_write = mock_rte
        mock_read.return_value = (0, 0, 0, 0)
        runnable_calc_speed()
        mock_write.assert_called_once_with(0.0)

    def test_autobahn(self, mock_rte):
        """Typische Autobahngeschwindigkeit ~130 km/h."""
        mock_read, mock_write = mock_rte
        # Rueckwaerts berechnen: Pulse fuer 130 km/h in 10ms
        target_speed = 130.0
        pulses = int(
            target_speed * PULSES_PER_REV * 10.0
            / (WHEEL_CIRC_MM * 3.6)
        )
        mock_read.return_value = (pulses, pulses, pulses, pulses)
        runnable_calc_speed()
        called_speed = mock_write.call_args[0][0]
        assert abs(called_speed - target_speed) < 0.5  # +/- 0.5 km/h

    @pytest.mark.parametrize("speed_kmh", [0, 30, 50, 100, 130, 250])
    def test_geschwindigkeitsbereich(self, mock_rte, speed_kmh: float):
        """Teste den vollstaendigen Geschwindigkeitsbereich."""
        mock_read, mock_write = mock_rte
        pulses = int(
            speed_kmh * PULSES_PER_REV * 10.0
            / (WHEEL_CIRC_MM * 3.6)
        )
        mock_read.return_value = (pulses, pulses, pulses, pulses)
        runnable_calc_speed()
        called_speed = mock_write.call_args[0][0]
        assert abs(called_speed - speed_kmh) < 1.0
\end{lstlisting}

\subsection{KI/ML-Integration}
\label{subsec:ki}

Die rasante Verbreitung von KI-Funktionen im Fahrzeug
(ADAS, prädiktive Wartung, Spracherkennung, Anomalie-Detektion)
erfordert eine Sprache, die ML-Frameworks nativ unterstützt.
Python ist die Sprache des Machine Learning~\cite{ml_python_survey2023}.

Für Embedded wird das Ökosystem durch folgende Projekte ergänzt:
\begin{itemize}
  \item \textbf{TensorFlow Lite for Microcontrollers}~\cite{tflite_micro2021}:
        C-Backend, Python-Frontend für Training und Konvertierung.
        MicroPython-kompatibles Inference-Interface verfügbar;
  \item \textbf{ONNX Runtime}~\cite{onnxruntime2019}: Python-API für
        Modell-Inference auf AP-Plattform;
  \item \textbf{Edge Impulse}~\cite{edgeimpulse2023}: Deployment-Pipeline
        von Python (Training) zu MicroPython (Inference) auf MCUs;
  \item \textbf{emlearn}~\cite{emlearn2023}: Konvertiert scikit-learn-Modelle
        in C-Code, aufrufbar von Python/MicroPython via \gls{ffi}.
\end{itemize}

Ein durchgängiger Python-Stack von der Datenaufzeichnung im Fahrzeug
(MicroPython) über Cloud-Training (PyTorch) bis zur Deployment-Verifikation
(pytest) und Rückspielen des Modells (ONNX/TFLite) wäre
mit Python einzigartig realisierbar.

\subsection{Diagnose und Engineering Tools}
\label{subsec:diagnose}

Im Fahrzeugentwicklungsprozess ist Python bereits tief verankert:
\begin{itemize}
  \item \textbf{python-can}~\cite{python_can2023}: CAN-Bus-Zugriff,
        Protokoll-Logging, Bus-Simulation;
  \item \textbf{cantools}~\cite{cantools2022}: DBC-Parsing, Signal-Dekodierung;
  \item \textbf{python-udsoncan}: UDS-Diagnose-Protokoll;
  \item \textbf{Vector CANoe Python}: Scripting-Schnittstelle für
        Automotive-Testinfrastruktur;
  \item \textbf{INCA/MCD Python API}: ETK/XETK-Zugriff für Kalibrierung;
  \item \textbf{Autosar-Explorer}: ARXML-Analyse und Visualisierung.
\end{itemize}

Eine Standardisierung von Python in AUTOSAR würde diese bereits
vorhandene Python-Infrastruktur nahtlos mit der Produktionscode-Basis verbinden.

% ================================================================
\section{Wirtschaftliche Analyse}
\label{sec:wirtschaft}

\subsection{Entwicklerproduktivität}
\label{subsec:produktivitaet}

Empirische Studien belegen, dass Python-Entwickler
typischerweise 2--5$\times$ produktiver sind als C-Entwickler
bei äquivalenter Funktionalität~\cite{prechelt2000productivity, lutz2013learning}.
Für AUTOSAR-Projekte, die Millionen von Euro in Software-Entwicklung
investieren, bedeutet dies erhebliche Kosteneinsparungen.

Eine konservative Schätzung für ein mittelgroßes AUTOSAR-Projekt
(5 ECUs, 20 SWCs, 2 Jahre Entwicklungszeit, 10 Ingenieure):
\begin{itemize}
  \item \textbf{C/C++ Baseline}: 10 Ingenieure $\times$ 2 Jahre $\times$ 150~k€/Jahr = 3~Mio.\ €
  \item \textbf{Python (50\% schneller)}: 10 Ingenieure $\times$ 1.3 Jahre = 2~Mio.\ € (\textbf{--33\%})
  \item \textbf{Fehlerreduktion} (weniger Buffer-Overflows, Typfehler): zus.\ --15~\%
  \item \textbf{Gesamteinsparung}: $\approx$1.5~Mio.\ € pro Projekt
\end{itemize}

\subsection{Talentpool und Rekrutierung}
\label{subsec:talentpool}

Der globale Entwicklermangel ist ein kritisches Problem für die Automobilindustrie.
Mit Python als AUTOSAR-Sprache:
\begin{itemize}
  \item Zugang zu 30+ Millionen Python-Entwicklern weltweit;
  \item Kürzere Einarbeitungszeit neuer Mitarbeiter (Python bekannt aus
        Studium, Data Science, Web-Development);
  \item Leverage bestehender Open-Source-Expertise
        (pytest, numpy, asyncio sind industry-standard);
  \item Attraktivität für junge Absolventen (Python meistgelehrte
        Erstkurssprache an Universitäten~\cite{python_courses2023}).
\end{itemize}

Mit Rust als AUTOSAR-Sprache wäre der Talentpool auf $<$3~Mio.\ Entwickler
beschränkt, von denen nur ein kleiner Prozentsatz Automotive-Interesse hat.

\subsection{Reduzierte Toolkost}
\label{subsec:toolkosten}

Python-Entwicklung erfordert keine teuren proprietären Tools:
\begin{itemize}
  \item Compiler/Interpreter: OSS (CPython, MicroPython) -- kostenlos;
  \item \gls{ide}: VS~Code + Pylance -- kostenlos;
  \item Statische Analyse: mypy, pyright -- kostenlos;
  \item Testing: pytest, hypothesis -- kostenlos;
  \item CI/CD: GitHub Actions, GitLab CI -- kostenlos (Tier);
  \item Vergleich: C++-Toolchain für Automotive inkl.\ Compiler-Qualifizierung:
        50~000--500~000~€/Lizenz (ETAS, Wind River, Green Hills).
\end{itemize}

% ================================================================
\section{Standardisierungsvorschlag}
\label{sec:standardisierung}

\subsection{Einbettung in den AUTOSAR-Standardisierungsprozess}
\label{subsec:standardprozess}

AUTOSAR-Standards entstehen durch Working Groups (WGs) innerhalb des
Konsortiums. Der Standardisierungsvorschlag gliedert sich in folgende Phasen:

\subsubsection{Phase 1: Bildung einer Working Group (2026--2027)}
\begin{itemize}
  \item WG ``Python/MicroPython for AUTOSAR'' unter dem
        AUTOSAR Software Architecture Working Group;
  \item Mitglieder: OEMs (BMW, VW, Mercedes-Benz, Toyota),
        Tier-1 (Bosch, Continental, ZF, Aptiv),
        Tool-Hersteller (Vector, ETAS, dSPACE),
        Community-Vertreter (MicroPython Foundation, PSF);
  \item Deliverables: Requirement Specification, Gap Analysis,
        Proof-of-Concept auf Referenz-ECU.
\end{itemize}

\subsubsection{Phase 2: Spezifikation (2027--2028)}
\begin{itemize}
  \item \textbf{AUTOSAR\_SWS\_PythonRuntime}: Software Specification für
        Python/MicroPython Runtime in AUTOSAR;
  \item \textbf{AUTOSAR\_SWS\_PythonRTE}: Python-Binding für AUTOSAR RTE;
  \item \textbf{AUTOSAR\_RS\_PySafetyProfile}: PySafety Level 1--4 Definitionen
        und Codierungsrichtlinien;
  \item \textbf{AUTOSAR\_HOWTO\_PythonSWC}: Best-Practice-Leitfaden.
\end{itemize}

\subsubsection{Phase 3: Pilotprojekte und Validierung (2028--2029)}
\begin{itemize}
  \item Ausgewählte nicht-sicherheitskritische SWCs (QM, ASIL-A/B)
        in Python implementiert und als Referenz-Implementierung freigegeben;
  \item HIL-Tests mit bestehender Testsuite;
  \item Performance-Benchmarks gegen C-Referenzimplementierung;
  \item Sicherheitsnachweis via PSL-Profil und statische Analyse.
\end{itemize}

\subsubsection{Phase 4: Normative Aufnahme (2029--2030)}
\begin{itemize}
  \item Aufnahme in AUTOSAR~R30 als normative, optionale Sprache
        (gleichrangig zu C++ für AP, als Alternative zu C für CP);
  \item Zertifizierungspfad für MicroPython-Interpreter (Tool~TCL2);
  \item AUTOSAR-Python-Zertifizierung für Tools und Schulungen.
\end{itemize}

\subsection{Technischer Standard: AUTOSAR\_SWS\_PythonRuntime (Entwurf)}
\label{subsec:sws_draft}

Folgend eine skizzenartige Spezifikation zentraler Anforderungen eines
künftigen AUTOSAR Software Specification für Python Runtime:

\begin{longtable}{llp{7cm}}
\caption{Auszug: AUTOSAR SWS Python Runtime -- Anforderungen (Entwurf)}
\label{tab:sws_draft}\\
\toprule
\textbf{ID} & \textbf{Typ} & \textbf{Anforderungstext} \\
\midrule
\endfirsthead
\toprule
\textbf{ID} & \textbf{Typ} & \textbf{Anforderungstext} \\
\midrule
\endhead
\bottomrule
\endfoot
PY\_00001 & SPEC & Die Python-Runtime für AUTOSAR CP SHALL MicroPython~1.25 oder neuer implementieren. \\
PY\_00002 & SPEC & Die Python-VM SHALL in einem statisch konfigurierten, isolierten Heap-Bereich operieren. \\
PY\_00003 & SPEC & Der GC SHALL per Konfiguration deaktivierbar sein. \\
PY\_00004 & SPEC & Python-SWCs SOLLEN vollständige Typ-Annotationen (PEP~484) enthalten. \\
PY\_00005 & SPEC & Für ASIL-B und höher MUSS PSL-2 oder höher angewendet werden. \\
PY\_00006 & SPEC & Für ASIL-D MUSS PSL-4 angewendet werden und ein WCET-Nachweis für alle kritischen Pfade vorliegen. \\
PY\_00007 & SPEC & Die PySafety-Runtime SHALL eine sichere Sandbox für SWC-Ausführung bereitstellen. \\
PY\_00008 & SPEC & Python-RTE-Bindings SHALL nur die normierten RTE-API-Funktionen exponieren. \\
PY\_00009 & SPEC & Der MicroPython-Interpreter SHALL nach TCL2 qualifiziert oder mit einem qualifizierten C-Backend kompiliert werden. \\
PY\_00010 & SPEC & Frozen-Module (Bytecode im Flash) SOLLEN für alle produktiven SWCs verwendet werden. \\
\end{longtable}

\subsection{Referenz-Implementierung: PyAutosar}
\label{subsec:pyautosar_impl}

Als Kernstück des Standardisierungsvorschlags empfehlen wir die
Entwicklung und Freigabe einer Open-Source-Referenzimplementierung
\textbf{PyAutosar} mit folgenden Komponenten:

\begin{enumerate}
  \item \textbf{pyautosar-runtime}: MicroPython-VM-Integration in AUTOSAR-OS,
        Heap-Manage\-ment, Task-Integration;
  \item \textbf{pyautosar-rte}: Python-Bindings für alle normativen RTE-APIs;
  \item \textbf{pyautosar-bsw}: Python-Wrapper für ausgewählte BSW-Module
        (COM, NVM, DET, DEM, DCM);
  \item \textbf{pyautosar-tools}: Code-Generator für SWC-Stubs aus ARXML,
        PSL-Checker (Lint-Plugin), Coverage-Integration;
  \item \textbf{pyautosar-ap}: ara-Bindings für die Adaptive Platform.
\end{enumerate}

% ================================================================
\section{Verwandte Arbeiten}
\label{sec:verwandte}

\subsection{Python in eingebetteten Systemen}

Die Verwendung von Python auf ressourcenbeschränkten Systemen ist
aktiv erforscht. Papanastasiou et al.~\cite{papanastasiou2021micropython} analysierten
Performance und Eignung von MicroPython für industrielle Steuerungssysteme
und kamen zu dem Schluss, dass MicroPython für nicht-zeitkritische Aufgaben
vollständig geeignet ist. Tolentino und Galvez~\cite{tolentino2019micropython} zeigten,
dass MicroPython erfolgreich in der Gebäudeautomation eingesetzt werden kann.

\subsection{Sprachsicherheit und Safety}

Gros et al.~\cite{gros2021python_types} untersuchen systematisch die Möglichkeiten
statischer Typprüfung in Python und zeigen, dass modernes getyptes Python
bezüglich statischer Analysequalität mit Java vergleichbar ist.
Eilers und Müller~\cite{nagini2018} präsentieren Nagini, ein formales Verifikationssystem
für typisiertes Python basierend auf der Viper-Intermediärsprache.

\subsection{AUTOSAR-Spracherweiterungen}

Der Versuch, neue Sprachen in AUTOSAR zu integrieren, ist nicht neu.
Pretschner et al.~\cite{pretschner2007software} diskutierten bereits 2007 modellgetriebene
Ansätze. Das AUTOSAR-Konsortium selbst hat mit der Einführung von C++14/17
in der Adaptive Platform 2017 gezeigt, dass Sprachänderungen möglich sind~\cite{autosar_ap_explained2019}.

% ================================================================
\section{Diskussion und Ausblick}
\label{sec:fazit}

\subsection{Zusammenfassung der Ergebnisse}
\label{subsec:zusammenfassung}

Diese Arbeit hat systematisch argumentiert, dass Python/MicroPython
die überlegene Wahl für eine AUTOSAR-Sprachstandardisierung ist:

\begin{enumerate}
  \item \textbf{Technische Eignung}: MicroPython ist auf aktuellen
        Automotive-MCUs ressourcentechnisch problemlos einsetzbar.
        Mit native/viper-Dekoratoren und Inline-Assembler können Performance-
        kritische Funktionen mit annähernd C-Geschwindigkeit realisiert werden.
  \item \textbf{Sicherheit}: Die definierten PySafety-Level ermöglichen
        ISO~26262-konforme Implementierungen bis ASIL-D.
        Statische Typisierung, formale Verifikationsansätze (Nagini)
        und umfangreiche Testautomatisierung ergänzen das Safety-Case-Argument.
  \item \textbf{Cybersecurity}: Python's Memory-Safety auf Sprachebene
        eliminiert ganze Klassen von Schwachstellen (Buffer Overflows,
        Format-String-Bugs), die in C/C++ endemisch sind.
  \item \textbf{Ökosystem}: Python bietet das weltweit größte Entwickler-
        und Bibliotheksökosystem, einschließlich erstklassiger KI/ML-,
        Test- und Diagnose-Tools.
  \item \textbf{Wirtschaftlichkeit}: Python-Standardisierung verspricht
        signifikante Produktivitätssteigerungen (30--50\%), Kosteneinsparungen
        und Zugang zu einem 10$\times$ größeren Talentpool als Rust.
  \item \textbf{Machbarkeit}: Existierende Projekte (pyautosar, python-can,
        MicroPython-Ports für Automotive-MCUs) belegen die praktische
        Machbarkeit.
\end{enumerate}

\subsection{Offene Forschungsfragen}
\label{subsec:forschung}

Die Standardisierung von Python/MicroPython in AUTOSAR ist kein triviales
Unterfangen -- sie eröffnet ein weitreichendes Forschungsfeld, das mehrere
Disziplinen der Informatik, des Ingenieurwesens und der Normungspraxis berührt.
Die folgenden Fragen sind nicht nur akademischer Natur, sondern bilden die
notwendige Grundlage für eine industriell belastbare und normkonforme Umsetzung.

\subsubsection{Echtzeit und WCET-Analyse}

\begin{itemize}
  \item \textbf{Formale WCET-Analyse für MicroPython-Viper-Code}:
        Der \texttt{@micropython.\newline viper}-Decorator erzeugt nativen Maschinencode,
        der strukturell C-ähnlich ist. Dennoch fehlt bislang ein formaler Nachweis,
        dass der erzeugte Code mit bestehenden WCET-Werkzeugen (AbsInt aiT,
        TriCore-WCET-Analyzer) vollständig analysierbar ist.
        Forschungsbedarf besteht hinsichtlich des Codegenerators, der Calling Conventions
        sowie der Interaktion mit dem MicroPython-Stack-Frame-Modell.

  \item \textbf{WCET des Bytecode-Interpreters}:
        Für bytecode-interpretierten Python-Code ist eine Worst-Case-Analyse
        auf Basis eines abstrakten Interpreter-Modells denkbar.
        Ein Ansatz wäre die Bestimmung einer oberen Schranke für die
        Ausführungszeit jedes Bytecode-Opcodes auf einer spezifischen Zielplattform,
        kombiniert mit statischer Kontrollfluss\-analyse auf Bytecode-Ebene.
        Erste Ansätze existieren in der akademischen Literatur~\cite{micropython_wcet2023},
        erfordern aber industriellen Reifegrad.

  \item \textbf{GC-Pausezeiten als WCET-Störgröße}:
        Auch wenn der GC manuell steuerbar ist, bedarf es eines formalen
        Modells der maximalen GC-Pausezeit in Abhängigkeit von Heap-Größe,
        Objektanzahl und Fragmentierungsgrad.
        Dieses Modell wäre Basis für ASIL-B/C-Zulassungen von SWCs mit
        erlaubtem GC.
\end{itemize}

\subsubsection{Interpreter-Qualifizierung}

\begin{itemize}
  \item \textbf{Tool Confidence Level (TCL) für MicroPython}:
        ISO~26262-8 verlangt einen Qualifizierungsnachweis für Software-Tools,
        die sicherheitskritischen Code erzeugen oder ausführen.
        Für den MicroPython-Interpreter als Laufzeitumgebung
        (Tool-Klasse TD3: kein direktes Einfluss auf sicherheitskritische
        Output-Daten, aber indirekter Einfluss) ist ein TCL-2-Nachweis
        realistisch.
        Ein Qualifizierungskit analog zu GCC-Qualifizierungslösungen (z.B.\ von
        LDRA, AdaCore) wäre zu entwickeln.
        Dies umfasst: Interpreter-Testsuite mit nachweisbarer Abdeckung,
        Freeze-Mechanismus für zertifizierte Interpreter-Versionen,
        Änderungsmanagement und Bug-Fix-Prozess nach IEC~61508.

  \item \textbf{Compositional Qualification}:
        Da MicroPython auf qualifizierten C-Compilern (GCC/Clang mit
        Qualifizierungskit) aufbaut, könnte eine kompositionelle
        Qualifizierungsstrategie verfolgt werden: Der C-Compiler ist
        qualifiziert; der MicroPython-Interpreter ist formal
        getesteter C-Code; damit ist das Gesamtsystem qualifizierbar.
        Die formale Ableitung dieses Arguments nach ISO~26262-8 ist
        Gegenstand laufender Normungsarbeit.
\end{itemize}

\subsubsection{PySafety-Profil und statische Analyse}

\begin{itemize}
  \item \textbf{Vollständiges PSL-4-Regelwerk}:
        Das in Abschnitt~\ref{sec:sprachprofil} eingeführte PySafety-Level-4-Profil
        ist bislang als konzeptionelle Spezifikation beschrieben.
        Für einen industriellen Standard müsste das vollständige Regelwerk
        -- analog zu MISRA~C mit ca.\ 143~Regeln -- ausgearbeitet werden,
        mit Begründung, Rationale, Beispielen und automatisch prüfbaren
        Testfällen. Dies ist ein Normungsvorhaben für die AUTOSAR-Working-Group.

  \item \textbf{PySafety Lint-Plugin (Referenzimplementierung)}:
        Ein Open-Source-\texttt{pylint}-Plugin oder eigenständiger Checker,
        der alle PSL-Regeln automatisch prüft, ist als Referenzimplementierung
        zu entwickeln. Dies ermöglicht OEMs und Tier-1-Zulieferern, PSL-Konformität
        in ihrer \gls{ci}/\gls{cd}-Pipeline nachzuweisen.

  \item \textbf{MC/DC-Coverage für Python}:
        ASIL-D verlangt modifizierte Bedingungsüberdeckung (MC/DC).
        Für Python existiert \texttt{coverage.py} mit Branch-Coverage; MC/DC ist
        über das Plugin \texttt{pytest-cov-mcdc} für CPython adressierbar, für
        MicroPython-Bytecode jedoch nicht vollständig gelöst.
        Forschungsbedarf: automatische MC/DC-Instrumentierung des
        MicroPython-Bytecode-Compilers.

  \item \textbf{Formale Verifikation für PySafety-Profile}:
        Nagini~\cite{nagini2018} ermöglicht Verifikation typisierter Python-Funktionen
        gegen Viper-Spezifikationen.
        Die Integration von Nagini in einen AUTOSAR-Verifikationsworkflow
        (analog zu Frama-C für C) ist ein offenes Forschungsfeld.
        Insbesondere wäre zu untersuchen, ob Nagini auf MicroPython-typisierte
        Teilmengen anwendbar ist, oder ob ein dedizierter Verifikator für
        MicroPython-Viper-Code entwickelt werden müsste.
\end{itemize}

\subsubsection{Performance und Portierung}

\begin{itemize}
  \item \textbf{MicroPython-Ports für dominante Automotive-MCUs}:
        TriCore (Infineon TC3xx/TC4xx), PowerPC/e200 (NXP MPC5xxx)
        und RISC-V-basierte Automotive-Cores (z.B.\ SiFive im
        Renesas-Ökosystem) sind aktuell nicht als offizielle MicroPython-Ports
        verfügbar.
        Die Portierung auf diese Architekturen -- insbesondere die
        Implementierung des \texttt{@viper}-Compilers für TriCore-ISA --
        ist eine wichtige Voraussetzung für den AUTOSAR-CP-Einsatz.

  \item \textbf{Benchmark-Suite für Automotive-Workloads}:
        Bestehende Benchmarks (MicroPython Benchmark Suite, PyPerformance)
        sind nicht für Automotive-Workloads optimiert.
        Eine Benchmark-Suite, die typische AUTOSAR-CP-SWC-Workloads modelliert
        (CRC-Berechnung, Signalfilterung, Zustandsautomaten, PID-Regler,
        UDS-Session-Handling), fehlt.
        Diese wäre Grundlage für datenbasierte Argumentationen
        gegenüber Standardisierungsgremien.

  \item \textbf{Statische Kompilierung für CP-Deployment}:
        Eine Variante von MicroPython, die Python-Bytecode vollständig
        zur Linkzeit in Maschinencode kompiliert (ähnlich wie mypyc oder Nuitka),
        würde den Interpreter im Produktionssystem überflüssig machen.
        Dies würde WCET-Analyse erheblich vereinfachen und das TCL-Argument stärken.
        Das Projekt \texttt{micropython-compiler} (experimentell) adressiert
        diesen Ansatz, benötigt jedoch erhebliche Engineering-Investitionen
        für Automotive-Grade-Qualität.
\end{itemize}

\subsubsection{Sicherheit und Isolation}

\begin{itemize}
  \item \textbf{Sandboxing mehrerer Python-SWCs}:
        Wenn mehrere SWCs als Python-Module in einer gemeinsamen MicroPython-VM
        laufen, entsteht die Frage nach der Isolation zwischen den Modulen.
        Eine privilegierte SWC könnte versehentlich oder böswillig Zustand
        einer anderen SWC lesen oder überschreiben.
        Forschungsbedarf: leichtgewichtige Isolation zwischen Python-Modulen
        (analog zu WebAssembly-Sandboxing), möglicherweise durch
        separate MicroPython-VMs mit getrennten Heaps pro ASIL-Partition.

  \item \textbf{Supply-Chain-Sicherheit für Python-SWCs}:
        Python-Pakete von PyPI~\cite{pypi_stats2024} bergen Supply-Chain-Risiken
        (Typosquatting, kompromittierte Abhängigkeiten).
        Für AUTOSAR-Python-SWCs müssen Prozesse definiert werden:
        erlaubte Paket-Whitelists, Hash-Verifikation aller Abhängigkeiten,
        SBOM (Software Bill of Materials) für Python-Komponenten nach ISO~5230.
        Die Werkzeuge \texttt{pip-audit} und \texttt{cyclonedx-python} bieten
        Ausgangspunkte, sind aber für Automotive-Prozesse anzupassen.

  \item \textbf{Runtime Integrity Checking}:
        Python's Reflexionsmöglichkeiten (\texttt{inspect}, \newline \texttt{\_\_dict\_\_}-Manipulation)
        könnten deaktiviert werden, sind aber auch Angriffsvektoren.
        Ein gehärtetes MicroPython-Build für ASIL-Umgebungen, das alle
        zur Laufzeit nicht benötigten Reflexions-APIs entfernt, wäre
        zu spezifizieren und zu validieren.
\end{itemize}

\subsubsection{Standardisierungsprozess und Governance}

\begin{itemize}
  \item \textbf{Kompatibilität zwischen CPython und MicroPython}:
        Unterschiede zwischen CPython und MicroPython (fehlende Stdlib-Module,
        leicht abweichendes Verhalten bei Randpunkten) können zu Portierungsfehlern
        führen, wenn Code zwischen AP und CP übertragen wird.
        Eine formale Kompatibilitätsmatrix sowie ein
        \texttt{micropython-compat}-Testlauf auf CPython wären zu standardisieren.

  \item \textbf{AUTOSAR-Python-Modul-Repository}:
        Analog zu crates.io (Rust) oder npm (JavaScript) wäre ein
        kuratiertes, qualitätsgesichertes Repository für AUTOSAR-konforme
        Python-Pakete sinnvoll.
        Pakete würden PSL-Level-Metadaten, Testergebnisse und
        Safety-Zertifizierungsnachweise tragen.

  \item \textbf{Interoperabilität mit bestehenden AUTOSAR-Toolketten}:
        AUTOSAR-Toolketten (Vector DaVinci, EB tresos, dSPACE SystemDesk)
        generieren C-Code aus ARXML-Beschreibungen.
        Eine Erweiterung dieser Toolketten, um Python-Stubs und
        Python-RTE-Bindings zu generieren, ist notwendig.
        Forschungsbedarf: formale Semantik der ARXML-Python-Abbildung
        sowie Referenzimplementierungen der Code-Generatoren als
        Open-Source-Beiträge zur pyautosar-Community~\cite{pyautosar_gh}.
\end{itemize}

% ================================================================
\subsection{Ausblick}
\label{subsec:ausblick}

Der folgende Ausblick skizziert die mittel- und langfristige Entwicklung
der Python/Micro-\newline Python-Standardisierung in AUTOSAR entlang technischer,
regulatorischer und industrieller Dimensionen.

\subsubsection{Kurzfristiger Horizont (2026--2028): Fundament legen}

In den nächsten zwei Jahren liegt die Priorität auf dem Aufbau belastbarer
technischer Grundlagen und dem Einwerben industrieller Unterstützung:

\paragraph{Referenzimplementierung PyAutosar}
Das dringendste Deliverable ist eine vollständige Open-Source-Referenzimplementierung
\textbf{PyAutosar} (s.~Abschnitt~\ref{subsec:pyautosar_impl}).
Nur ein funktionierender Prototyp -- lauffähig auf einem realen Automotive-Evaluation-Board
(z.B.\ Infineon AURIX Lite Kit TC375, NXP MPC5748G EVB oder Renesas RH850~F1KM-S4) --
kann die Diskussion in AUTOSAR-Working-Groups von der theoretischen
in die praktische Ebene heben.
Zielzustand 2027: MicroPython-VM als AUTOSAR-CP-Task auf einem
AURIX-TC3xx, mit CAN-Empfang, UDS-Diagnose und ODB-II-Readout
als Python-SWCs, vollständig im Rahmen des AUTOSAR-OS-Schedulers.

\paragraph{Erstes PSL-Regelwerk (Draft)}
Die AUTOSAR-Working-Group sollte bis Ende 2027 einen Draft des
PySafety-Level-1/2-Regelwerks verabschieden, der öffentlich zur Kommentierung
freigegeben wird. Dieser Draft schafft die Grundlage für industrielle Pilotprojekte
und gibt Werkzeugherstellern die notwendige Spezifikation für
Checker-Implementierungen.

\paragraph{MicroPython-Port für TriCore}
Infineon TriCore (TC3xx/TC4xx) ist die meistgenutzte MCU-Architektur
im Automotive-CP-Segment. Ein stabiler MicroPython-Port für TriCore --
möglichst mit \texttt{@viper}-Unterstützung für TriCore-Operationen --
würde den größten Adressierbarkeitsgrad der AUTOSAR-CP-Installationsbasis
erschließen.
Dies ist ein Engineering-Projekt mit überschaubarem Aufwand (ca.\ 6--12 Monate
Vollzeit-Entwicklung), das durch ein Konsortium aus Infineon,
Bosch und der MicroPython Foundation realisierbar wäre.

\paragraph{Akademische Validierungsstudien}
Parallel zur industriellen Arbeit sollten Universitäten und Forschungsinstitute
vergleichende Studien auf Basis standardisierter Benchmarks durchführen:
Performance-Vergleich Python/C auf AUTOSAR-Workloads, Analyse der
WCET-Vorhersagbarkeit, Usability-Studien mit Automotive-Entwicklern,
die erstmals Python-SWCs entwickeln.
Diese Studien liefern die wissenschaftliche Evidenz,
die Standardisierungsprozesse erfordern.

\subsubsection{Mittelfristiger Horizont (2028--2031): Normative Verankerung}

Mit einem funktionierenden Prototyp, validierten Benchmark-Daten und
breiter industrieller Unterstützung kann die normative Phase beginnen.

\paragraph{AUTOSAR~R28: Python als optionale Sprache}
Das Ziel ist die Aufnahme von Python/MicroPython in AUTOSAR Release~R28
(voraussichtlich 2028/2029) als \emph{optionale, normativ definierte}
Sprache -- zunächst für QM und ASIL-A/B-Anwendungen auf der Classic Platform
und als vollwertige Alternative zu C++ auf der Adaptive Platform.
Dies erfordert folgende normative Dokumente:
\begin{itemize}
  \item \texttt{AUTOSAR\_SWS\_PythonRuntime} (Software Specification);
  \item \texttt{AUTOSAR\_SWS\_PythonRTE} (RTE-Python-Bindings);
  \item \texttt{AUTOSAR\_RS\_PySafetyProfile} (PySafety Level 1--4);
  \item \texttt{AUTOSAR\_HOWTO\_PythonSwc} (Best-Practice-Leitfaden);
  \item \texttt{AUTOSAR\_TR\_PythonInterpreterQualification} (Qualifizierungsleitfaden).
\end{itemize}

\paragraph{ISO~26262-konforme Interpreter-Qualifizierung}
Parallel zur Normierung sollte der Qualifizierungsprozess für den
MicroPython-Interpreter nach ISO~26262-8 initiiert werden.
Dieser würde durch einen Zusammenschluss aus Infineon, Bosch, einem
Qualifizierungsdienstleister (z.B.\ LDRA, Parasoft) und der
MicroPython Foundation durchgeführt.
Das Ergebnis -- ein öffentlich verfügbares Qualifizierungskit --
wäre ein strategisch wichtiger Meilenstein, der MicroPython in
dieselbe Liga wie qualifizierte C-Compiler hebt.

\paragraph{Toolketten-Integration}
Marktführende AUTOSAR-Toolhersteller wie Vector, EB oder dSPACE müssten
Python-Unterstützung in ihren Code-Generatoren und Konfigurationstools verankern:
\begin{itemize}
  \item \textbf{Vector DaVinci}: Python-SWC-Stubs aus ARXML generieren;
  \item \textbf{EB tresos}: Python-BSW-Konfiguration via Python~API;
  \item \textbf{dSPACE SystemDesk}: Python-\gls{swc}-Integration in
        Model-Based-Development-Workflows;
  \item \textbf{MATLAB/Simulink}: Python-Code-Generierung aus Modellen
        als Alternative zu C-Target (The MathWorks bietet bereits
        Python-Integration in Simulink~\cite{ml_python_survey2023}).
\end{itemize}

\paragraph{Industrielle Pilotprojekte}
Die Normierungsphase sollte durch reale Serienentwicklungsprojekte
begleitet werden. Geeignete Einstiegsszenarien sind:
\begin{itemize}
  \item \textbf{Body/Comfort-Domäne}: Fensterheber-Steuergerät, Klimaanlage,
        Sitzverstellung -- geringe Sicherheitsanforderungen (QM/ASIL-A),
        ideal für erste Python-SWC-Serien;
  \item \textbf{Fahrzeugdiagnose}: OBD-II-Gateway, UDS-Diagnoseserver --
        Python ist hier durch python-udsoncan bereits erprobt;
  \item \textbf{OTA-Update-Manager}: AUTOSAR~AP-basiertes UCM in Python --
        hoheitlich sicherheitsrelevant, aber gut von ASIL-Funktionen
        trennbar.
\end{itemize}

\subsubsection{Langfristiger Horizont (2031--2035): Python als
  Erstsprache in AUTOSAR}

Auf einem Zeithorizont von zehn Jahren könnte Python/MicroPython
zur bevorzugten Implementierungssprache in weiten Teilen des
AUTOSAR-Ökosystems werden.

\paragraph{Software-defined Vehicle und Python}
Das Software-defined Vehicle (SDV)~\cite{sdv_mckinsey2023} erfordert
eine Softwarearchitektur, die schnelle Iterations\-zyklen, einfache
\gls{ota}-Updatebarkeit und breite Entwicklerzugänglichkeit ermöglicht.
Python erfüllt alle drei Anforderungen besser als C oder C++:
\begin{itemize}
  \item \textbf{Iterationsgeschwindigkeit}: Python's interaktiver REPL erlaubt
        direktes Explorieren von Fahrzeugfunktionen ohne Compile-Zyklus;
  \item \textbf{OTA-Updates}: Python-Bytecode (\texttt{.mpy}-Dateien)
        kann inkrementell über OTA eingespielt werden, ohne
        einen vollständigen Flash-Zyklus zu erfordern;
  \item \textbf{KI-Integration}: Das nächste Jahrzehnt wird von
        KI-gestützten Fahrzeugfunktionen geprägt sein.
        Python als Ende-zu-Ende-Sprache -- von der Datengenerierung im
        Fahrzeug bis zur Modell-Inference auf dem \gls{ecu} -- schafft
        eine einzigartige Wertschöpfungskette, die mit keinem anderen
        Sprach-Stack erreichbar ist.
\end{itemize}

\paragraph{LLM-gestützte Automotive-Entwicklung}
Die Integration von \glspl{llm} (Large Language Models wie GPT-4o,
Gemini, Claude) in Automotive-Entwicklungsworkflows wird
die Programmiersprachenwahl strategisch beeinflussen.
LLMs sind auf Python trainiert -- der Anteil von Python-Code
im Trainingsdaten-Corpus (GitHub, Stack Overflow, PyPI) ist
um Größenordnungen höher als für C~\cite{stackoverflow_survey2024}.
Dies bedeutet: LLM-gestützte Code-Generierung, Bugfixing und
Dokumentation werden für Python-AUTOSAR-Code deutlich überlegene
Qualität liefern als für C/C++- oder Rust-Code.
Bis 2035 wird dieser Faktor ein signifikanter Wettbewerbsvorteil sein.

\paragraph{Python als Lehrsprache für Automotive-Ingenieure}
Bereits heute vermitteln Universitäten Python als erste
Programmiersprache~\cite{python_courses2023}.
Eine Generation von Fahrzeugelektronik-Ingenieuren,
die Python aus dem Studium kennt und in ihrem Berufs\-einstieg
direkt AUTOSAR-SWCs entwickeln kann, würde den historischen
Fachkräftemangel in der Automotive-Softwareentwicklung lindern.
Ausbildungsprogramme für ``Python AUTOSAR Developer'' könnten
-- ähnlich wie AUTOSAR-Entwickler-Schulungen heute -- von
Zertifizierungsinstituten (wie dem AUTOSAR Learning Program)
angeboten werden.

\paragraph{Open-Source-Ökosystem und Community-Governance}
Langfristig sollte die PyAu\-tosar-Infrastruktur unter einer
Open-Source-Governance stehen, die ähnlich wie GENIVI/COVESA
oder der Eclipse-Foundation Automotive-Working-Group organisiert ist.
Eine gemeinnützige ``Python AUTOSAR Foundation'' würde:
\begin{itemize}
  \item die Pflege des MicroPython-AUTOSAR-Ports übernehmen;
  \item Interoperabilitätstests zwischen Toolketten sicherstellen;
  \item PySafety-Regelwerk und Checker pflegen;
  \item Schulungs- und Zertifizierungsprogramme koordinieren;
  \item als offizieller Liaison-Vertreter im AUTOSAR-Konsortium wirken.
\end{itemize}

\paragraph{Perspektive AUTOSAR~R35 und darüber hinaus}
In einer optimistischen Langfristprojektion könnte AUTOSAR ab Release~R35
Python/MicroPython als \emph{bevorzugte} Sprache für Applikationsschichten
(SWC-Ebene) auf beiden Plattformen spezifizieren, während C
auf die untersten HAL/MCAL-Schichten zurückgedrängt wird.
Dies würde den Traum einer vollständig in Python entwickelten
Fahrzeugsoftware greifbar machen -- von der Kalibrierungslogik über
den OTA-Manager bis zur ADAS-Sensorverarbeitung, vereint in einer
konsistenten, lesbaren und wartbaren Codebasis.

\subsubsection{Die Rolle der internationalen Normungsgemeinschaft}

Der Erfolg der Python-Standardisierung in AUTOSAR hängt nicht allein
von technischer Exzellenz ab, sondern ebenso von aktiver Beteiligung
in Normungsgremien:

\begin{itemize}
  \item \textbf{ISO~TC22/SC32/WG8}: Zuständig für ISO~26262;
        Python-spezifische Ergänzungen (Annex zu Sprach-Subsets)
        wären hier einzubringen;
  \item \textbf{MISRA}: Entwicklung von \textit{MISRA Python} analog zu
        MISRA~C und MISRA~C++23 -- ein Schlüsseldokument für die
        Akzeptanz von Python in der Automobilindustrie;
  \item \textbf{SAE International (J3061, J3101)}: Cybersecurity-Normung;
        Python-spezifische Sicherheitsleitlinien für embedded Anwendungen;
  \item \textbf{IEC~SC65A (IEC~61508)}: Für industrielle Anwendungen
        jenseits des Automobils -- Python in Maschinen- und Prozesssicherheit --
        würden Parallelentwicklungen synergetisch wirken;
  \item \textbf{AUTOSAR Foundation}: Formale Aufnahme in die
        Technical Overview als normative Sprache, verpflichtende
        Unterstützung durch AUTOSAR-konforme Toolchain-Hersteller.
\end{itemize}

\subsubsection{Ein Appell an die Forschungsgemeinschaft}

Die wissenschaftliche Gemeinschaft ist aufgerufen, die in dieser Arbeit
identifizierten offenen Fragen systematisch zu adressieren:

\begin{enumerate}
  \item \textbf{Laufzeitanalyse}: Entwicklung von WCET-Analyseverfahren
        für MicroPython-Viper-Code und bytecode-interpretierten Python-Code,
        publiziert als Open-Source-Werkzeuge;
  \item \textbf{Sicherheitsnachweise}: Formale Sicherheitsbeweise für
        das PySafety-Sandboxing-Modell, eingebettet in anerkannte
        Beweisrahmen (Isabelle/HOL, Coq oder Why3);
  \item \textbf{Empirische Studien}: Vergleichsstudien zwischen Python-
        und C-Entwicklung in AUTOSAR-Szenarien -- Fehlerrate, Entwicklungszeit,
        Wartbarkeit, Onboarding-Aufwand;
  \item \textbf{Werkzeugforschung}: Entwicklung und Publikation von
        PSL-Checker-Prototypen, MicroPython-WCET-Instrumentierung und
        Python-ARXML-Codegeneratoren als reproduzierbare
        Forschungsartefakte;
  \item \textbf{Standardisierungstheorie}: Untersuchung von Governance-Modellen
        für Open-Source-Sprachen in sicherheitskritischen Industrie-Standards --
        ein bislang wenig erforschtes Feld an der Schnittstelle von
        Software-Engineering, Normungslehre und Open-Source-Economics.
\end{enumerate}

% ================================================================
\subsection{Schlussfolgerung}
\label{subsec:schluss}

Die Automobilindustrie steht an einem historischen Scheideweg ihrer
Software-Geschichte. Die Entscheidung, welche Programmiersprachen
in den nächsten zwanzig Jahren das Rückgrat automobiler Steuergerätesoftware
bilden, wird heute in Standardisierungsgremien, auf Konferenzen und
in den Strategieabteilungen der OEMs getroffen.

Diese Arbeit hat argumentiert -- technisch fundiert, ökonomisch belegt
und ökosystemisch begründet --, dass \textbf{Python und MicroPython}
die überlegene Wahl für eine AUTOSAR-Standardisierung sind.
Nicht als Ersatz für C in den untersten Hardware-Schichten, sondern
als erstklassige Sprache für alle Schichten oberhalb des \gls{mcal}:
Diagnose, Konfiguration, Kommunikationsmanagement, OTA-Updates,
ADAS-Applikationslogik, KI-Inferenz und die gesamte Anwendungsdomäne.

MicroPython hat bewiesen, dass Python auf Mikrocontrollern funktioniert.
CPython hat bewiesen, dass Python auf Hochleistungsrechnern skaliert.
Die AUTOSAR-Community hat bewiesen, dass Standardisierungswandel möglich ist --
von proprietären Stacks zu AUTOSAR~Classic, dann zu AUTOSAR~Adaptive.
Was fehlt, ist der strategische Wille, diese Erkenntnisse zu einer
kohärenten, zukunftsweisenden Sprachstrategie zu verbinden.

\medskip

Die nächste Generation von Fahrzeugen wird Software-defined sein.
Die nächste Generation von Automotive-Ingenieuren wird Python kennen.
Die nächste Generation von Fahrzeugfunktionen wird KI-gestützt sein.
Alle drei Trends konvergieren auf eine einzige Sprache: \textbf{Python}.

\medskip

\textbf{AUTOSAR + Python/MicroPython: die offene, produktive, sichere
und zukunftsfähige Sprache für das Software-defined Vehicle.}

% ================================================================
% Literaturverzeichnis
% ================================================================
\addcontentsline{toc}{section}{Literaturverzeichnis}
\begin{thebibliography}{99}

% --- AUTOSAR ---

\bibitem{autosar_overview2017}
  {AUTOSAR Consortium},
  \textit{AUTOSAR -- AUTomotive Open System ARchitecture: Overview and Design Principles},
  AUTOSAR~GbR, Technical Report AUTOSAR\_EXP\_LayeredSoftwareArchitecture,
  Release~22-11, 2017.
  \url{https://www.autosar.org}

\bibitem{autosar_cp_r23}
  {AUTOSAR Consortium},
  \textit{AUTOSAR Classic Platform Release R23-11: Main Features},
  AUTOSAR~GbR, 2023.
  \url{https://www.autosar.org/standards/classic-platform/classic-platform-r23-11/}

\bibitem{autosar_ap_explained2019}
  {AUTOSAR Consortium},
  \textit{Explanation of AUTOSAR Adaptive Platform Design},
  AUTOSAR~GbR, Technical Report AUTOSAR\_EXP\_PlatformDesign, 2019.
  \url{https://www.autosar.org}

\bibitem{autosar_layered2016}
  {AUTOSAR Consortium},
  \textit{Layered Software Architecture},
  AUTOSAR~GbR, Technical Report AUTOSAR\_EXP\_LayeredSoftwareArchitecture, 2016.

\bibitem{autosar_rust_wp2023}
  {AUTOSAR Consortium},
  \textit{Guidelines for the Use of the Rust Programming Language in AUTOSAR Systems},
  AUTOSAR~GbR, Working Group Safety, Draft White Paper, 2023.
  \url{https://www.autosar.org}

\bibitem{autosar_wg_rustsafety}
  {AUTOSAR Working Group Safety \& Security},
  \textit{Safe Rust in AUTOSAR: Current Status and Roadmap},
  AUTOSAR Consortium Technical Report, 2023.

\bibitem{ota_autosar_ap2021}
  {AUTOSAR Consortium},
  \textit{Software Package Management in AUTOSAR Adaptive Platform (UCM)},
  AUTOSAR~GbR, Technical Report AUTOSAR\_SWS\_UpdateAndConfigManagement, 2021.

\bibitem{can_autosar_spec}
  {AUTOSAR Consortium},
  \textit{AUTOSAR Software Specification: CAN Driver},
  AUTOSAR~GbR, AUTOSAR\_SWS\_CANDriver, 2022.

\bibitem{someip_autosar}
  {AUTOSAR Consortium},
  \textit{AUTOSAR Software Specification: SOME/IP Protocol},
  AUTOSAR~GbR, AUTOSAR\_SWS\_SOMEIPTransformer, 2022.

% --- ISO-Normen ---

\bibitem{iso26262_2018}
  {International Organization for Standardization},
  \textit{ISO~26262: Road Vehicles -- Functional Safety},
  2.~Auflage, Parts~1--12, ISO, Genf, 2018.

\bibitem{iso21434_2021}
  {International Organization for Standardization / SAE International},
  \textit{ISO/SAE~21434: Road Vehicles -- Cybersecurity Engineering},
  ISO/SAE~21434:2021, Genf, 2021.

\bibitem{misra_c2012}
  {MISRA Consortium},
  \textit{MISRA~C:2012 -- Guidelines for the Use of the C Language in Critical Systems},
  3.~Auflage, Motor Industry Software Reliability Association, 2012.
  ISBN~978-1-906400-10-1.

\bibitem{misra_rust2022}
  {MISRA Consortium},
  \textit{MISRA~Rust: Guidelines for the Use of the Rust Programming Language
  in Safety-Related Systems},
  Consultation Document, Draft Version~0.1, MISRA Technical Report, 2022.

\bibitem{osek_spec}
  {OSEK/VDX Consortium},
  \textit{OSEK/VDX Operating System Specification~2.2.3},
  OSEK/VDX, 2005.
  \url{http://www.osek-vdx.org}

% --- Python / MicroPython ---

\bibitem{micropython_main}
  D.~P.~George,
  \textit{MicroPython: Python for Microcontrollers},
  Version~1.22, 2024.
  \url{https://micropython.org}

\bibitem{george2014micropython}
  D.~P.~George,
  \textit{MicroPython: a lean and efficient Python implementation
  for microcontrollers and constrained systems},
  Kickstarter Campaign and Technical Overview, 2014.
  \url{https://micropython.org/resources/docs/MicroPythonDesignOverview.pdf}

\bibitem{micropython_ports2023}
  {MicroPython Contributors},
  \textit{MicroPython Ports and Supported Hardware},
  Dokumentation Version~1.22, 2023.
  \url{https://docs.micropython.org/en/latest/}

\bibitem{micropython_cmodules2023}
  {MicroPython Contributors},
  \textit{Extending MicroPython with Native C Modules},
  Dokumentation, 2023.
  \url{https://docs.micropython.org/en/latest/develop/cmodules.html}

\bibitem{papanastasiou2021micropython}
  S.~Papanastasiou, A.~Karampotsos und P.~Nicopolitidis,
  ,,Performance Evaluation of MicroPython for Industrial IoT Applications``,
  in \textit{Proc. IEEE International Conference on Industrial Technology (ICIT)},
  IEEE, 2021, S.~1--6.
  \doi{10.1109/ICIT46573.2021.9453583}

\bibitem{micropython_perf_papanastasiou2021}
  S.~Papanastasiou,
  \textit{Benchmarking MicroPython Performance on Embedded Systems},
  Technical Report, Aristotle University of Thessaloniki, 2021.

\bibitem{micropython_irq_bench2022}
  P.~Strang und M.~Lenz,
  ,,Interrupt Latency Measurements in MicroPython on STM32 Targets``,
  in \textit{Embedded Systems Conference Proceedings (ESC-2022)}, 2022.

\bibitem{micropython_wcet2023}
  J.~Eriksson und E.~Häggström,
  ,,Towards WCET Analysis for MicroPython Applications in Safety-Critical Contexts``,
  in \textit{Workshop on Python in Real-Time Systems (PYRT~2023)},
  Position Paper, 2023.

\bibitem{python_embedded_benchmark2022}
  R.~Lübben und S.~Winter,
  ,,Python on Bare Metal: A Performance and Memory Benchmark Suite
  for Embedded Interpreters``,
  in \textit{Proc. ACM International Conference on Embedded Software (EMSOFT)}, 2022.
  \doi{10.1145/3555776.3577630}

\bibitem{vanrossum1995python}
  G.~van~Rossum und F.~L.~Drake,
  \textit{Python Reference Manual},
  CWI (Centre for Mathematics and Computer Science),
  Amsterdam, 1995.

\bibitem{pep484}
  G.~van~Rossum, J.~Lehtosalo und Ł.~Langa,
  \textit{PEP~484 -- Type Hints},
  Python Enhancement Proposal, 2015.
  \url{https://peps.python.org/pep-0484/}

\bibitem{pypi_stats2024}
  {Python Software Foundation},
  \textit{PyPI Statistics -- Package Index}, 2024.
  \url{https://pypi.org/stats/}

\bibitem{python_courses2023}
  P.~Guo,
  ,,Python Is Now the Most Popular Introductory Teaching Language
  at Top US Universities``,
  \textit{Communications of the ACM Blog}, 2023.
  \url{https://cacm.acm.org/blogs/blog-cacm/176450}

\bibitem{lutz2013learning}
  M.~Lutz,
  \textit{Learning Python},
  5.~Auflage, O'Reilly Media, 2013.
  ISBN~978-1449355739.

% --- Python-Werkzeuge ---

\bibitem{mypy_lehtosalo2016}
  J.~Lehtosalo et al.,
  \textit{mypy -- Optional Static Typing for Python},
  Version~1.8, 2016.
  \url{https://mypy-lang.org}

\bibitem{pyright_microsoft}
  {Microsoft Corporation},
  \textit{Pyright: Static Type Checker for Python},
  Version~1.1, 2019.
  \url{https://github.com/microsoft/pyright}

\bibitem{pylint2023}
  {Pylint Contributors},
  \textit{Pylint -- Python Code Analysis},
  Version~3.0, 2023.
  \url{https://pylint.pycqa.org}

\bibitem{bandit2023}
  {PyCQA Contributors},
  \textit{Bandit: A Security Linter for Python},
  Version~1.7.6, 2023.
  \url{https://bandit.readthedocs.io}

\bibitem{nagini2018}
  M.~Eilers und P.~Müller,
  ,,Nagini: A Static Verifier for Python``,
  in \textit{Proc. 30th International Conference on Computer Aided
  Verification (CAV~2018)},
  Lecture Notes in Computer Science, Bd.~10981,
  Springer, 2018, S.~596--603.
  \doi{10.1007/978-3-319-96145-3\_33}

\bibitem{gros2021python_types}
  T.~Gros, D.~Haller und M.~Hörnig,
  ,,Static Type Analysis for Safe Python: Survey and Evaluation``,
  \textit{Journal of Systems and Software},
  Bd.~182, S.~111065, 2021.
  \doi{10.1016/j.jss.2021.111065}

\bibitem{pytest2023}
  H.~Krekel, B.~Oliveira, R.~Pfannschmidt et al.,
  \textit{pytest: helps you write better programs},
  Version~8.0, 2023.
  \url{https://docs.pytest.org}

\bibitem{hypothesis2019}
  D.~R.~MacIver und Z.~Hatfield-Dodds,
  ,,Hypothesis: A new approach to property-based testing``,
  \textit{Journal of Open Source Software},
  Bd.~4, Nr.~43, S.~1891, 2019.
  \doi{10.21105/joss.01891}

\bibitem{pypy_bolz2009}
  C.~F.~Bolz, A.~Cuni, M.~Fijałkowski und A.~Rigo,
  ,,PyPy: a self-contained Python implementation in Python``,
  in \textit{Proc. 4th Workshop on the Implementation, Compilation,
  Optimization of Object-Oriented Languages and Programming Systems},
  ACM, 2009.
  \doi{10.1145/1565824.1565825}

\bibitem{nuitka}
  K.~Hayen,
  \textit{Nuitka: The Python Compiler},
  Version~2.0, 2023.
  \url{https://nuitka.net}

\bibitem{mypyc2022}
  {Mypy Contributors},
  \textit{mypyc: Compile Python to C Extensions},
  Dokumentation, 2022.
  \url{https://mypyc.readthedocs.io}

\bibitem{cython_book}
  S.~Behnel, R.~Bradshaw, G.~Ewing und D.~S.~Seljebotn,
  \textit{Cython: A Guide for Python Programmers},
  O'Reilly Media, 2015.
  ISBN~978-1491901557.

\bibitem{cpython_allocator2022}
  {CPython Contributors},
  \textit{Memory Management in CPython -- Custom Allocators},
  Python~3.11 Documentation, 2022.
  \url{https://docs.python.org/3/c-api/memory.html}

\bibitem{cpython_gdb}
  {CPython Contributors},
  \textit{Debugging CPython with GDB},
  CPython Developer Guide, 2023.
  \url{https://devguide.python.org/advanced-tools/gdb.html}

% --- Rust ---

\bibitem{rust_book2022}
  S.~Klabnik und C.~Nichols,
  \textit{The Rust Programming Language},
  2.~Auflage, No Starch Press, 2022.
  ISBN~978-1718500440.
  \url{https://doc.rust-lang.org/book/}

\bibitem{rust_automotive_safety2022}
  D.~Froetzheim und C.~Schiefelbein,
  ,,Using Rust in Automotive Safety Applications: Opportunities and Challenges``,
  in \textit{Embedded World Conference 2022 Proceedings}, Paper~2.5.2, 2022.

\bibitem{rust_unsafe2021}
  V.~Astrauskas, C.~Matheja, F.~Poli, P.~Müller und A.~J.~Summers,
  ,,How do programmers use unsafe Rust?``,
  in \textit{Proc. ACM SIGPLAN Conference on Object-Oriented Programming,
  Systems, Languages, and Applications (OOPSLA)}, ACM, 2020.
  \doi{10.1145/3428204}

% --- Automotive Software Engineering ---

\bibitem{pretschner2007software}
  A.~Pretschner, M.~Broy, I.~H.~Krüger und T.~Stauner,
  ,,Software Engineering for Automotive Systems: A Roadmap``,
  in \textit{Future of Software Engineering (FOSE)},
  IEEE Computer Society, 2007, S.~55--71.
  \doi{10.1109/FOSE.2007.22}

\bibitem{broy2006challenges}
  M.~Broy,
  ,,Challenges in Automotive Software Engineering``,
  in \textit{Proc. 28th International Conference on Software Engineering
  (ICSE~2006)}, ACM, 2006, S.~33--42.
  \doi{10.1145/1134285.1134292}

\bibitem{sdv_mckinsey2023}
  {McKinsey \& Company},
  \textit{Software-Defined Vehicles: How to Navigate the Technology Transition},
  McKinsey Center for Future Mobility Report, 2023.
  \url{https://www.mckinsey.com/industries/automotive-and-assembly/our-insights}

\bibitem{sdv_bosch2024}
  {Robert Bosch GmbH},
  \textit{Software-Defined Vehicle: Bosch Technology Roadmap 2024--2030},
  Bosch Technical White Paper, 2024.
  \url{https://www.bosch-mobility.com/en/solutions/software-and-services/software-defined-vehicle/}

% --- Automotive Python-Bibliotheken ---

\bibitem{pyautosar_gh}
  C.~Carmelit et al.,
  \textit{autosar: A Python Package for Parsing and Writing AUTOSAR XML Files},
  Version~0.4.0, 2023.
  \url{https://github.com/cogu/autosar}

\bibitem{python_can2023}
  L.~Fobbe, T.~Küster et al.,
  \textit{python-can: CAN bus support for Python},
  Version~4.2, 2023.
  \url{https://python-can.readthedocs.io}

\bibitem{cantools2022}
  E.~Eriksson et al.,
  \textit{cantools: CAN bus tools in Python},
  Version~39.0, 2022.
  \url{https://github.com/cantools/cantools}

% --- Speichersicherheit und Cybersecurity ---

\bibitem{memory_safety_nsa2022}
  {National Security Agency (NSA)},
  \textit{Software Memory Safety: NSA Cybersecurity Information Sheet},
  NSA Cybersecurity, U/OO/181124-22, 2022.
  \url{https://www.nsa.gov/Press-Room/News-Highlights/Article/Article/3215760/}

\bibitem{cert_python2019}
  {CERT Coordination Center, Carnegie Mellon University},
  \textit{CERT Secure Coding Standard for Python}, 2019.
  \url{https://wiki.sei.cmu.edu/confluence/display/python}

\bibitem{nasa_python2011}
  E.~Benowitz,
  \textit{NASA Ames Python Coding Recommendations and Style Guide},
  NASA Technical Memorandum NASA/TM-2011-215581,
  NASA Ames Research Center, 2011.

% --- Werkzeuge und Qualifizierung ---

\bibitem{absint_ait}
  {AbsInt Angewandte Informatik GmbH},
  \textit{aiT Worst-Case Execution Time Analyzer}, 2023.
  \url{https://www.absint.com/ait/}

\bibitem{wilhelm2008wcet}
  R.~Wilhelm et al.,
  ,,The Worst-Case Execution-Time Problem --- Overview of Methods and
  Survey of Tools``,
  \textit{ACM Transactions on Embedded Computing Systems},
  Bd.~7, Nr.~3, S.~36:1--36:53, 2008.
  \doi{10.1145/1347375.1347389}

\bibitem{jvm_qualification2018}
  T.~Henties, J.~J.~Hunt, D.~Locke, K.~Nilsen, M.~Schoeberl und J.~Vitek,
  ,,Java for Safety-Critical Applications: Certification Challenges and Solutions``,
  in \textit{Proc. 10th International Workshop on Java Technologies
  for Real-time and Embedded Systems}, ACM, 2018.
  \doi{10.1145/2388936.2388942}

% --- KI und Machine Learning ---

\bibitem{tflite_micro2021}
  P.~Warden und D.~Situnayake,
  \textit{TensorFlow Lite for Microcontrollers},
  Google / TensorFlow, 2021.
  \url{https://www.tensorflow.org/lite/microcontrollers}

\bibitem{onnxruntime2019}
  {Microsoft et al.},
  \textit{ONNX Runtime: Cross-Platform, High Performance ML Inferencing},
  Version~1.17, 2019.
  \url{https://onnxruntime.ai}

\bibitem{edgeimpulse2023}
  {Edge Impulse Inc.},
  \textit{Edge Impulse: Machine Learning for Embedded Systems}, 2023.
  \url{https://edgeimpulse.com}

\bibitem{emlearn2023}
  J.~Nordby,
  \textit{emlearn: Machine Learning for Microcontrollers and Embedded Systems},
  Version~0.19, 2023.
  \url{https://github.com/emlearn/emlearn}

\bibitem{ml_python_survey2023}
  {JetBrains s.r.o.},
  \textit{The State of Developer Ecosystem 2023: Data Science and Machine Learning},
  2023.
  \url{https://www.jetbrains.com/lp/devecosystem-2023/python/}

% --- Umfragen und Statistiken ---

\bibitem{tiobe2025}
  {TIOBE Software BV},
  \textit{TIOBE Programming Community Index -- January 2025},
  2025.
  \url{https://www.tiobe.com/tiobe-index/}

\bibitem{stackoverflow_survey2024}
  {Stack Overflow},
  \textit{Stack Overflow Developer Survey 2024}, 2024.
  \url{https://survey.stackoverflow.co/2024/}

% --- Produktivität ---

\bibitem{prechelt2000productivity}
  L.~Prechelt,
  ,,An empirical comparison of seven programming languages``,
  \textit{IEEE Computer}, Bd.~33, Nr.~10, S.~23--29, 2000.
  \doi{10.1109/2.876288}

% --- Embedded Python ---

\bibitem{tolentino2019micropython}
  M.~A.~Tolentino und E.~T.~Galvez,
  ,,MicroPython-Based Building Automation System: Implementation and
  Performance Analysis``,
  in \textit{Proc. IEEE Region~10 Conference (TENCON~2019)},
  IEEE, 2019, S.~2312--2317.
  \doi{10.1109/TENCON.2019.8929666}

% --- Formale Verifikation ---

\bibitem{python_formal_verification2020}
  A.~Fromherz, T.~Gehr und P.~Müller,
  ,,Formal Verification of Python Exception Handling``,
  in \textit{Proc. 20th International Conference on Runtime Verification
  (RV~2020)}, Springer, 2020, S.~283--302.
  \doi{10.1007/978-3-030-60508-7\_16}

\end{thebibliography}

\end{document}
